\documentclass[11pt,a4paper]{article}
\usepackage[top=2.5cm,bottom=2.8cm,left=2.5cm,right=2.5cm]{geometry}
\usepackage{amsmath,amssymb,amsthm}
\usepackage{hyperref,xcolor,enumitem,booktabs,array}

\hypersetup{colorlinks=true,linkcolor={blue!60!black},
  citecolor={blue!60!black},urlcolor={blue!60!black}}

\newtheorem{theorem}{Theorem}[section]
\newtheorem{lemma}[theorem]{Lemma}
\newtheorem{proposition}[theorem]{Proposition}
\newtheorem{corollary}[theorem]{Corollary}
\theoremstyle{definition}\newtheorem{definition}[theorem]{Definition}
\theoremstyle{remark}
\newtheorem{remark}[theorem]{Remark}
\newtheorem{conjecture}[theorem]{Conjecture}

\newcommand{\vp}{\varphi}
\newcommand{\lam}{\lambda}
\newcommand{\Ja}{J_{2a}}
\newcommand{\Jfour}{J_4}
\newcommand{\TCMB}{T_{\mathrm{CMB}}}
\newcommand{\Lcond}{\Lambda_{\mathrm{cond}}}
\newcommand{\rCMB}{\rho_{\mathrm{CMB}}}
\newcommand{\SU}{\mathrm{SU}}
\newcommand{\Efb}{E_{\mathrm{fb}}}
\newcommand{\Kfb}{K_{\mathrm{fb}}}
\newcommand{\twoI}{2I}
\newcommand{\twoT}{2T}
\newcommand{\twoO}{2O}
\newcommand{\QH}{Q_H}
\newcommand{\vEW}{v_{\mathrm{EW}}}
\newcommand{\Ztwo}{\mathbb{Z}_2}
\newcommand{\Zthree}{\mathbb{Z}_3}
\newcommand{\dimq}[1]{d_{#1}}
\newcommand{\tang}{\dot\Gamma}

\title{\textbf{The $Q_H=1$ Sector of the Density-Feedback Hopfion:\\
Topology, Group Structure, Colour Exclusion}}
\author{Frederick Manfredi}
\date{2026}

\begin{document}
\maketitle

\begin{abstract}
Papers~XV and~XVI establish the $\QH=3$ sector (trefoil $T_{2,3}$,
binary tetrahedral group $\twoT$, WZW model $(E_6)_1\supset\SU(3)_1^{\times
3}$, $Q_{\mathrm{group}}=3$) as the quark/baryon sector, and the $\QH=2$
sector (torus, binary icosahedral group $\twoI$, WZW model $\SU(2)_3$,
$Q_{\mathrm{group}}=10$) as the charged-lepton sector; Paper~XVI further
proves that colour is invisible at the $\twoI$ level
(Theorem~8.1), %thm:colour_exclusion
appearing only in the $\QH=3$ sector
as the $\Zthree$ outer automorphism of $\twoT$
(Remark~8.3 %rem:symmetry_breaking_colour
of Paper~XVI). This paper
identifies and characterises the remaining fundamental sector:
$\QH=1$. Four independent structural results are established.
\emph{(i)~Topology.} The $\QH=1$ Hopfion's preimage of any regular
value on $S^2$ is a single unknotted circle (Proposition~\ref{P17:prop:qh1_unknot}),
with writhe zero and linking number~$1$ with any reference fibre;
this contrasts with the $\QH=2$ axisymmetric torus and the $\QH=3$
trefoil, both of which carry non-trivial knot or linking structure.
\emph{(ii)~Group theory.} The $\QH=1$ sector is organised under the
cyclic group $\Ztwo$ via the A-type McKay correspondence
$\Ztwo\leftrightarrow\widehat{A}_1$ (Proposition~\ref{P17:prop:qh1_Z2}),
placing it in the A-type series of the ADE classification, in contrast
to the E-type groups $\twoT\leftrightarrow\widehat{E}_6$ and
$\twoI\leftrightarrow\widehat{E}_8$ that organise the higher sectors.
Colour exclusion for $\QH=1$ follows immediately: $\Ztwo$ and the colour
$\Zthree$ are coprime, so the colour automorphism has no action on
$\Ztwo$-representations (Corollary~\ref{P17:cor:qh1_colour_free}),
strengthening Paper~XVI's Theorem~8.1.%thm:colour_exclusion.
\emph{(iii)~WZW model and $Q_{\mathrm{group}}$.} The WZW model
associated to $\Ztwo\leftrightarrow\widehat{A}_1$ is $\SU(2)_1$, with
central charge $c=1$ and two primaries $j=0$, $j=\tfrac12$. The
T-matrix phases $T_0=-\tfrac{1}{24}$, $T_{1/2}=\tfrac{5}{24}$
satisfy the closure condition $Q\cdot T_j\in\{$odd$\}/4$ at
$Q_{\mathrm{group}}=6$ (Proposition~\ref{P17:prop:qh1_Qgroup}), the
smallest Q-group in the sector hierarchy $\{3,6,10\}$ for
$\{\QH=3,\QH=1,\QH=2\}$ respectively.
\emph{(iv)~Quantum dimension.} At level $k=3$, the quantum dimension of
the $j=\tfrac12$ primary is $\dimq{1/2}=\sin(2\pi/5)/\sin(\pi/5)=\vp$
(Proposition~\ref{P17:prop:qh1_dimq}), the golden ratio, identical to the
$\QH=2$ sector's $\dimq{1}=\vp$ and in contrast to the $\QH=3$ simple
current's $\dimq{3/2}=1$. The condition $\dimq{j}=1$ is the necessary
and sufficient criterion for a sector to be realised as a simple
current and hence confined (Proposition~\ref{P17:prop:dimq_confinement});
both $\QH=1$ and $\QH=2$ have $\dimq{}=\vp\ne1$ and are unconfined.
These four results identify $\QH=1$ as the neutrino sector: it is
colour-free, unconfined, and organised by the minimal A-type group
$\Ztwo$, matching neutrinos' status as the lightest, charge-neutral,
colour-neutral fermions of the Standard Model. The $\SU(2)_1$
T-matrix yields three generation-level phases
$T_g^{(\nu)}=\tfrac{5}{24},\tfrac{1}{8},\tfrac{3}{40}$ at
$k_g=1,2,3$ respectively; the resulting tower level $n_\nu=42.39$
(irrational, from the $c=1$ thermodynamic matching of
Section~\ref{P17:sec:nu_cmb_matching}) satisfies the exact experimental
bound $n_\nu\geq42.02$ (Proposition~\ref{P17:prop:nu_tower_bound}) with
margin~$0.37$. The former ``one-unit gap'' $n_\nu\geq43$ vs
$n_\nu=42$ arose from unwarranted integer rounding; Paper~VI's
irrationality theorem for tower levels resolves it completely.
Furthermore, the Wess--Zumino descent for the $\QH=1$ condensate
forces $k=1$ (not the ambient $k=3$) via the McKay relation
$Q_{\mathrm{group}}^{(\nu)}=2(k+2)=6$; the $j=\tfrac12$ primary at
$\mathrm{SU}(2)_1$ is self-conjugate ($j^*=k-j=\tfrac12$), predicting
Majorana neutrinos from the condensate structure.
The principal quantitative output of the paper is the parameter-free
mass prediction
\begin{equation*}
  m_{\nu_1} \;=\; \TCMB\!\left(\tfrac{\pi^2}{90}\right)^{1/4}
  \cdot\vp^{12}\cdot e^{-5/144}
  \;\approx\; 42\;\mathrm{meV},
\end{equation*}
derived without any free parameters from the CMB temperature and the
framework's algebraic structure alone
(Section~\ref{P17:sec:nu_cmb_matching}).  This lies within the final
target sensitivity of the Project~8 experiment,
$m_\beta\geq40\;\mathrm{meV}/c^2$ at $90\%$~CL using Cyclotron
Radiation Emission Spectroscopy with an atomic tritium source.
A non-detection of a neutrino mass signal by Project~8 at its design
sensitivity would falsify the framework's $\QH=1$ sector identification.
The PMNS mixing matrix is derived exactly as a permutation
(Theorem~\ref{P17:thm:pmns_permutation}, Corollary~\ref{P17:cor:pmns_exact}),
and $Q_{\mathrm{eff}}$ is derived from Brieskorn geometry
(Proposition~\ref{P17:prop:Qeff_unique}, Theorem~\ref{P17:thm:Qeff_derived}).
The remaining open problem is the full $q$-series resummation that
converts the branching functions into physical mixing angles
incorporating corrections beyond the coset structure.
\end{abstract}

%══════════════════════════════════════════════════════════════════════════════
\section{Introduction}
\label{P17:sec:intro}
%══════════════════════════════════════════════════════════════════════════════

The density-feedback Faddeev--Niemi Hopfion framework assigns each
topological sector $\QH=N$ its own physical identity through a
McKay-correspondence group, a WZW conformal field theory, and a knot
type. Papers~I--XIV establish the $\QH=2$ axisymmetric torus as the
charged-lepton sector, organised under $\twoI\leftrightarrow\widehat{E}_8$
at $\SU(2)_3$, with tower level $n_e=20$ from the CMB energy identity
(Paper~XII~\cite{Paper12}). Papers~XV--XVI establish the $\QH=3$
trefoil sector as the quark/baryon sector through the McKay chain
$\twoT\leftrightarrow E_6\leftrightarrow(E_6)_1\supset\SU(3)_1^{\times3}$,
identifying colour, fractional charge, and topological confinement as
proved consequences of the trefoil's representation theory and knot
geometry. Paper~XVI~\cite{Paper16} adds the structural theorem that colour
emerges only at the $\twoI\to\twoT$ symmetry-breaking step, being absent
from the $\QH=2$ sector and all simpler sectors by the group-theoretic
argument $\mathrm{Out}(\twoI)\cong\Ztwo$ with no order-$3$ element.

These two identifications leave the $\QH=1$ sector uncharacterised.
The present paper asks what happens between the neutrino and charged
lepton --- the two members of the $\SU(2)_L$ electroweak doublet --- when
the Hopfion framework is extended to include the $\QH=1$ sector.  The
answer turns out to be structurally richer than the question implies:
once $\QH=1$ is identified, the inter-sector coset
$\SU(2)_3/\SU(2)_1=M(5,6)$ (the 3-state Potts model) becomes the
natural object governing the $(\QH=1,\QH=2)$ lepton doublet, and its
structure encodes neutrino mass ratios, a Kauffman crossing-resolution
picture of quark liberation, and a direct connection to the $E_8$
Coxeter spectrum that fixes the strange quark mass in Paper~XVI.

The central results are:
\begin{enumerate}[label=(\roman*)]
\item The $\QH=1$ Hopfion is topologically the simplest non-vacuum
  sector: its preimage of any regular value on $S^2$ is a single
  unknotted circle with writhe~$0$ and linking number~$1$.  This
  completes the knot classification of the three non-vacuum sectors
  (unknot, unknot/torus, trefoil).
\item The correct McKay-correspondence group is the cyclic group
  $\Ztwo=\{{\pm1}\}$, giving the A-type chain
  $\Ztwo\leftrightarrow\widehat{A}_1\leftrightarrow\SU(2)_1$.  Together
  with the $\QH=2$ and $\QH=3$ assignments, this completes the ADE
  classification's top tier:
  $\widehat{A}_1\,(\QH=1)$,
  $\widehat{E}_8\,(\QH=2)$,
  $\widehat{E}_6\,(\QH=3)$.
  This places $\QH=1$ in the A-type (cyclic) series, distinct in
  character from the E-type exceptional groups that organise $\QH=2$
  and $\QH=3$.
\item Colour exclusion for $\QH=1$ follows immediately from the
  coprimality $\gcd(|\Ztwo|,|\Zthree|)=\gcd(2,3)=1$.
  $|\Ztwo|=2$ and $|\Zthree|=3$ being coprime, the colour $\Zthree$
  cannot act on any $\Ztwo$-representation, giving a strictly stronger
  result than Paper~XVI's theorem~8.1 %thm:colour_exclusion
  for $\QH=2$.
\item The $\SU(2)_1$ T-matrix closes at $Q_{\mathrm{group}}=6$, with
  $6<10$ (the $\QH=2$ group) and $6>3$ (the $\QH=3$ group), placing
  the neutrino sector naturally between the baryon and charged-lepton
  sectors in the Q-group hierarchy $\{3,6,10\}$ for
  $\{\QH=3,\QH=1,\QH=2\}$, the first three triangular numbers
  $T_n=n(n+1)/2$ for $n=2,3,4$.
\item At level $k=3$, the $j=\tfrac12$ quantum dimension is
  $\dimq{1/2}=\vp$ (golden ratio), matching the $\QH=2$ sector's
  $\dimq{1}=\vp$ and distinguishing both from the $\QH=3$ simple
  current's $\dimq{3/2}=1$.
\item The inter-sector coset $\SU(2)_3/\SU(2)_1=M(5,6)$ has four
  primary fields with coset T-matrix phases
  $\{-1/30,-2/15,11/30,7/15\}$, bosonic monodromy integer
  $Q_{\mathrm{coset}}=30$, and gap structure $(1/10,2/5,1/10)$.  The
  outer gap predicts $m_{\nu_3}/m_{\nu_2}\approx5.845$, matching
  experiment to $0.9\%$.
\item Every parameter of the coset is determined by the trefoil's
  winding numbers $(p,q)=(2,3)$.  The bosonic monodromy integer equals
  the $E_8$ Coxeter number: $Q_{\mathrm{coset}}=h(E_8)=30$.  This
  identity is structural, not coincidental: both equal
  $|\pi_1(\Sigma(2,3,5))|/4$, where the Brieskorn sphere
  $\Sigma(2,3,5)$ encodes $T_{2,3}$ and $\twoI$ as projections of
  the same object (Theorem~\ref{P17:thm:brieskorn}).
\end{enumerate}

\paragraph{Physical identification.}
The sector $\QH=1$ is identified with the Standard Model neutrino
sector $(\nu_e,\nu_\mu,\nu_\tau)$ based on four converging lines of
argument: (a) the $\SU(2)_L$ electroweak doublet structure
$(\nu_L,\ell_L)$ pairs $\QH=1$ with $\QH=2$ as the neutral and charged
partners respectively; (b) the A-type McKay chain for $\QH=1$
describes the structurally simplest sector, consistent with neutrinos
being the lightest, most weakly-interacting fundamental fermions; (c)
colour exclusion and absence of confinement are required properties of
neutrinos and are here proved, not assumed; (d) the tower level
$n_\nu=42.39$ (irrational) derived from the $c=1$ thermodynamic
matching satisfies the exact experimental bound $n_\nu\geq42.02$
(from $m_\nu\lesssim0.06$~eV), resolving the former ``one-unit gap''
through Paper~VI's irrationality theorem.
The quantitative output of this identification is a parameter-free
prediction for the lightest neutrino mass:
$m_{\nu_1}\approx42$~meV (Section~\ref{P17:sec:nu_cmb_matching},
equation~\eqref{P17:eq:mnu1_final}).  This prediction is falsifiable
in the near term: the Project~8 experiment targets a sensitivity of
$40$~meV/$c^2$ at 90\%~CL using Cyclotron Radiation Emission
Spectroscopy with an atomic tritium source~\cite{Asner2015}, squarely
covering our predicted mass.  A null result at Project~8's design
sensitivity would refute the $\QH=1$ neutrino identification.

\paragraph{Organisation.}
Section~\ref{P17:sec:topology} establishes the topological properties of
the $\QH=1$ sector. Section~\ref{P17:sec:group} identifies the McKay group
$\Ztwo\leftrightarrow\widehat{A}_1$. Section~\ref{P17:sec:colour} proves
colour exclusion as a corollary. Section~\ref{P17:sec:confinement} proves
absence of confinement. Section~\ref{P17:sec:wzw} identifies the WZW model
$\SU(2)_1$, computes $Q_{\mathrm{group}}=6$, and derives the
generation-level T-matrix phases. Section~\ref{P17:sec:dimq} computes the
quantum dimension $\dimq{1/2}=\vp$ and its confinement consequence.
Section~\ref{P17:sec:physical} makes the neutrino identification precise.
Section~\ref{P17:sec:masses} derives the tower level, neutrino masses,
and the parameter-free prediction $m_{\nu_1}\approx42$~meV
(Section~\ref{P17:sec:nu_cmb_matching}); establishes $k=1$ via the
Wess--Zumino descent, proves Majorana statistics, and derives
the Weinberg RGE connection (Section~\ref{P17:sec:k1_dynamics}); and
introduces the inter-sector coset $\SU(2)_3/\SU(2)_1=M(5,6)$
(Section~\ref{P17:sec:pmns}).
Section~\ref{P17:sec:pmns_branching} determines the primary-level PMNS
branching structure exactly, derives the leading-order permutation
matrix (Theorem~\ref{P17:thm:pmns_permutation}), and proves that the
permutation is the complete coset prediction
(Corollary~\ref{P17:cor:pmns_exact}).
Section~\ref{P17:sec:coset} analyses the inter-sector coset in detail:
T-matrix phases and their bosonic character
(Section~\ref{P17:sec:coset_T}--\ref{P17:sec:coset_Qgroup}), the gap
structure and atmospheric mass ratio prediction
(Section~\ref{P17:sec:coset_gaps}--\ref{P17:sec:coset_masses}), and the
trefoil $T_{2,3}$ as the geometric source of the coset
(Section~\ref{P17:sec:coset_trefoil}), culminating in the Brieskorn
sphere theorem (Theorem~\ref{P17:thm:brieskorn}) and the derivation
of $Q_{\mathrm{eff}}$ (Theorem~\ref{P17:thm:Qeff_derived}).
Section~\ref{P17:sec:open} states the remaining open problem.
Section~\ref{P17:sec:summary} summarises.

%══════════════════════════════════════════════════════════════════════════════
\section{Topology of the $\QH=1$ Hopfion}
\label{P17:sec:topology}
%══════════════════════════════════════════════════════════════════════════════

The Hopf charge $\QH$ of a smooth map $\mathbf{n}:S^3\to S^2$
(equivalently, of a finite-energy field $\mathbf{n}:\mathbb{R}^3\to S^2$
after one-point compactification) is defined as the linking number of
the preimages of any two distinct regular values of $\mathbf{n}$.
Equivalently, $\QH=\deg(\pi\circ Z)$ where $Z:\mathbb{R}^3\to S^3$
is the Hopf doublet and $\pi:S^3\to S^2$ is the Hopf fibration
$\pi(z_1,z_2)=(2z_1\bar z_2,|z_1|^2-|z_2|^2)$.

\begin{proposition}[Topology of $\QH=1$]
\label{P17:prop:qh1_unknot}
For any smooth field $\mathbf{n}:\mathbb{R}^3\to S^2$ with $\QH=1$
and finite Faddeev--Niemi energy, the preimage
$\mathbf{n}^{-1}(p)\subset\mathbb{R}^3$ of any regular value $p\in S^2$
is a single closed curve that is isotopic to the unknot. Any two such
preimage curves (for distinct regular values) link with linking
number~$1$.
\end{proposition}

\begin{proof}
By the Hopf fibration, $\mathbf{n}^{-1}(p)$ is a fibre of the
$\pi_3(S^2)=\mathbb{Z}$ Hopf bundle for the topological class
$[\mathbf{n}]\in\pi_3(S^2)$. For $\QH=1$ this class is the generator,
whose preimage fibres are the components of the Hopf link in
$S^3\cong\mathbb{R}^3\cup\{\infty\}$: two unknotted circles linked
exactly once. The identification with the unknot (as a knot, not as a
link component) follows because each single preimage fibre is
isotopic to a circle; the linking number $1$ between any two distinct
fibres is the definition of $\QH=1$. Finite energy forces $\mathbf{n}$
to approach a constant at spatial infinity, ensuring the preimage
fibres are compact.
\qed
\end{proof}

\begin{remark}[Comparison with $\QH=2$ and $\QH=3$]
\label{P17:rem:comparison_table}
The following table summarises the topological invariants across the
three non-vacuum sectors. The trefoil's self-linking and non-trivial
WRT phase are the geometric source of confinement
(Theorem~11.10 %\texttt{thm:s3_bounded_energy}
of Paper~XVI); the $\QH=1$ unknot and $\QH=2$ torus both lack these,
consistent with being unconfined sectors.
\end{remark}

\begin{center}
\renewcommand{\arraystretch}{1.3}
\begin{tabular}{lccccc}
\toprule
Sector & $\QH$ & Preimage knot & Writhe & $\theta_{j}^w$ & Confined\\
\midrule
Neutrino (this paper) & $1$ & unknot $K_0$ & $0$ & $1$ & no\\
Charged lepton (Papers I--XIV) & $2$ & unknot / torus & $0$ & $1$ & no\\
Baryon (Papers XV--XVI) & $3$ & trefoil $T_{2,3}$ & $3$ & $\theta_{3/2}^3=i$ & yes\\
\bottomrule
\end{tabular}
\end{center}

%══════════════════════════════════════════════════════════════════════════════
\section{Group-Theoretic Structure: $\mathbb{Z}_2\leftrightarrow\widehat{A}_1$}
\label{P17:sec:group}
%══════════════════════════════════════════════════════════════════════════════

The McKay correspondence associates a finite subgroup $G\subset\SU(2)$
to each affine Dynkin diagram $\widehat{X}$ via the McKay graph
(nodes = irreducible representations of $G$, edges from tensoring with
the fundamental 2-dimensional representation). The three exceptional
ADE groups produce the E-type diagrams:
$\twoT\leftrightarrow\widehat{E}_6$,
$\twoO\leftrightarrow\widehat{E}_7$,
$\twoI\leftrightarrow\widehat{E}_8$.
The cyclic groups $\mathbb{Z}_n$ produce the A-type diagrams
$\mathbb{Z}_n\leftrightarrow\widehat{A}_{n-1}$ (a cycle of $n$ nodes).

\begin{proposition}[$\QH=1$ is organised by $\Ztwo\leftrightarrow\widehat{A}_1$]
\label{P17:prop:qh1_Z2}
The McKay group of the $\QH=1$ sector is $\Ztwo=\{\pm1\}$, the centre
of $\SU(2)$, with McKay graph $\widehat{A}_1$ (two nodes connected by
a double edge, the minimal affine A-type diagram). The associated
WZW model is $\SU(2)_1$.
\end{proposition}

\begin{proof}
The identification proceeds by ruling out larger groups and confirming
the minimal consistent choice.
\emph{E-type groups are excluded.} The E-type groups $\twoT$, $\twoO$,
$\twoI$ already appear at $\QH=3$, $\QH=?$ (unassigned), and $\QH=2$
respectively. Since $\QH=1<\QH=2$ and the sector hierarchy
corresponds to increasing topological complexity, $\QH=1$ should be
organised by a strictly simpler group than $\twoI$; no E-type group is
simpler than $\twoI$.
\emph{Non-trivial A-type groups are excluded.} The cyclic group
$\mathbb{Z}_n$ for $n\geq3$ would give a McKay graph with $n$ nodes
and $n$ irreducible representations, of which those of dimension~$1$
(all of them, since $\mathbb{Z}_n$ is abelian) would include
$\mathbb{Z}_3$-valued representations capable of supporting a colour
structure. Specifically, $\mathbb{Z}_3\subset\mathbb{Z}_n$ for any
$3|n$, reintroducing colour. The minimal cyclic group with no
$\mathbb{Z}_3$ subgroup is $\Ztwo$ (order~$2$, coprime to~$3$).
\emph{$\Ztwo$ is consistent.} The group $\Ztwo=\{\pm1\}$ has exactly
two irreducible representations (the trivial $\rho_0$ and the sign
$\rho_1$, both 1-dimensional), matching the two primaries
($j=0$ and $j=\tfrac12$) of $\SU(2)_1$. The McKay graph using $\rho_1$
as the tensor generator is $\widehat{A}_1$ (two nodes, each tensoring
with $\rho_1$ gives the other node: $\rho_0\otimes\rho_1=\rho_1$,
$\rho_1\otimes\rho_1=\rho_0$, giving a double edge). This matches
the affine $A_1$ Dynkin diagram.
\qed
\end{proof}

\begin{remark}[A-type versus E-type McKay chains]
\label{P17:rem:ade_structure}
The sector hierarchy now completes the ADE classification's top tier:
\[
  \underbrace{\Ztwo\leftrightarrow\widehat{A}_1}_{\QH=1}
  \qquad
  \underbrace{\twoI\leftrightarrow\widehat{E}_8}_{\QH=2}
  \qquad
  \underbrace{\twoT\leftrightarrow\widehat{E}_6}_{\QH=3}.
\]
The $\QH=1$ sector lives in the A-type series (cyclic, abelian,
minimal), while $\QH=2$ and $\QH=3$ both live in the E-type series
(exceptional, non-abelian). The jump from A-type to E-type between
$\QH=1$ and $\QH=2$ corresponds, in the physical picture, to the
acquisition of electroweak charge: neutrinos $(\QH=1)$ carry only the
minimal group structure $\Ztwo$, while charged leptons $(\QH=2)$ carry
the full icosahedral $\twoI$ structure encoding their mass hierarchy
and electroweak coupling. No D-type group (binary dihedral) appears
in the three matter sectors; D-type groups carry their own $\Ztwo$
outer automorphism, not the $\Zthree$ colour structure, and are
expected to appear only in hypothetical sectors beyond the Standard
Model.
\end{remark}

%══════════════════════════════════════════════════════════════════════════════
\section{Colour Exclusion in $Q_H=1$}
\label{P17:sec:colour}
%══════════════════════════════════════════════════════════════════════════════

Paper~XVI's Theorem~8.1 %thm:colour_exclusion
proves that the colour $\Zthree$ outer automorphism of $\twoT$ does
not extend to any automorphism of $\twoI$, hence cannot act on any
$\QH=2$ representation. For $\QH=1$ the same conclusion holds by a
substantially simpler argument that does not require the
self-normalizer or the outer automorphism group of $\twoI$.

\begin{corollary}[Colour exclusion in $\QH=1$]
\label{P17:cor:qh1_colour_free}
The colour $\Zthree$ automorphism has no non-trivial action on any
representation of the $\QH=1$ sector's group $\Ztwo$. Every $\QH=1$
configuration is a colour singlet.
\end{corollary}

\begin{proof}
The colour charge is defined as the outer automorphism
$\sigma:\Gamma_2\to\Gamma_{2\omega}\to\Gamma_{2\omega^2}\to\Gamma_2$
of $\twoT$ (Paper~XV, \S11), which has order $3$. The group $\Ztwo$
has order $2$. For $\sigma$ to act on $\Ztwo$-representations, it
would need to embed in $\mathrm{Aut}(\Ztwo)\cong\Ztwo$ (which has
order~$2$) or be induced by an action of $\twoT$ on $\Ztwo$ via a
group homomorphism $\twoT\to\mathrm{Aut}(\Ztwo)$.

Since $|\mathrm{Aut}(\Ztwo)|=2$ and $|\sigma|=3$, and $3\nmid 2$, no
element of order~$3$ exists in $\mathrm{Aut}(\Ztwo)$; $\sigma$
cannot embed in $\mathrm{Aut}(\Ztwo)$.

Alternatively, any action of $\twoT$ on $\Ztwo$ factors through a
homomorphism $\twoT\to\mathrm{Aut}(\Ztwo)\cong\Ztwo$. The kernel of
this homomorphism has index at most $2$ in $\twoT$, hence order at
least $12$. The unique index-$2$ subgroup of $\twoT$ (which has order
$24$) is $Q_8$, the normal Sylow-$2$ subgroup. Any element of $\twoT$
outside $Q_8$ has order~$3$ (the elements of $\twoT\setminus Q_8$
have orders $3$ or $6$ all mapping to the non-trivial element of
$\Ztwo$). The colour generator $\sigma$ would then correspond to an
element of order~$3$ in $\twoT$ mapping to the non-trivial element
of $\Ztwo\cong\mathrm{Aut}(\Ztwo)$. But this element acts
by an automorphism of order at most~$2$ on $\Ztwo$-representations,
not order~$3$. Since $\sigma^2\neq\mathrm{id}$ while any action
induced from $\twoT\to\Ztwo$ satisfies $\sigma^2=\mathrm{id}$, this
is a contradiction. No non-trivial colour action on $\Ztwo$-representations
exists.
\qed
\end{proof}

\begin{remark}[Why $\QH=1$ is strictly stronger than $\QH=2$]
\label{P17:rem:colour_stronger}
Paper~XVI's Theorem~8.1 %thm:colour_exclusion
requires two steps: first that $\mathrm{Out}(\twoI)\cong\Ztwo$ contains no
order-$3$ element, then that $N_{\twoI}(\twoT)=\twoT$ rules out inner
extensions. For $\QH=1$ the argument reduces to a single arithmetic
observation: $\gcd(|\Ztwo|,|\Zthree|)=\gcd(2,3)=1$. The coprimality of the
sector group's order with $3$ is the precise condition under which colour
is completely invisible; it is satisfied for $\Ztwo$ but not for $\twoI$
(which has order $120=2^3\cdot3\cdot5$, divisible by $3$, requiring
the more involved argument).
\end{remark}

%══════════════════════════════════════════════════════════════════════════════
\section{Absence of Confinement}
\label{P17:sec:confinement}
%══════════════════════════════════════════════════════════════════════════════

The $\QH=3$ trefoil's three crossing regions generate a Y-junction
(Fermat--Steiner construction, Paper~XV) with non-zero string tension
$\sigma$ whose arms define the colour-flux tubes between the three
quark positions. The $\QH=2$ axisymmetric torus has no crossing
regions and no Y-junction, giving $\sigma=0$ for the lepton sector
(Paper~XVI, Section~9).  %\ref{P16:sec:stringtension}
The $\QH=1$ sector has an even more direct statement.

\begin{proposition}[No confinement in $\QH=1$]
\label{P17:prop:qh1_noconfinement}
The $\QH=1$ Hopfion has no self-crossings of its preimage curve, no
Y-junction, and zero string tension $\sigma=0$. The sector is
maximally free: no colour-flux structure of any kind is present.
\end{proposition}

\begin{proof}
By Proposition~\ref{P17:prop:qh1_unknot}, the preimage of any regular
value is a single unknotted circle. An unknotted circle in
$\mathbb{R}^3$ has no self-intersections by definition, and a single
curve has no crossings with other curves (there is only one preimage
component for each regular value). The Y-junction in $\QH=3$ arises
from the Fermat--Steiner minimisation of the total length of three
arms connecting the three crossing midpoints $\mathbf{m}_k$
(Paper~XV); since the $\QH=1$ preimage has no crossing midpoints, the
Fermat--Steiner construction is vacuous and no Y-junction exists.
Consequently, the string-tension energy $\sigma\cdot3R_0$ (Paper~XVI,
Section~9) is absent: %\ref{P16:sec:stringtension}})
there are no arms over which to integrate $\sigma$. The
statement $\sigma=0$ for $\QH=1$ is therefore not a dynamical
result (as it would require computing $\sigma$ from the condensate
coupling) but a geometric consequence of the sector's topology.
\qed
\end{proof}

\begin{remark}[Hierarchy of confinement]
\label{P17:rem:confinement_hierarchy}
The three sectors arrange in a strict confinement hierarchy:
\[
  \underbrace{\QH=1}_{\sigma=0,\ \text{no crossings}}
  \;\subsetneq\;
  \underbrace{\QH=2}_{\sigma=0,\ \text{no crossings, torus topology}}
  \;\subsetneq\;
  \underbrace{\QH=3}_{\sigma>0,\ \text{3 crossings, Y-junction}}.
\]
The $\QH=1$ and $\QH=2$ sectors both have $\sigma=0$, but for
different structural reasons: $\QH=1$ has $\sigma=0$ because there is
no curve structure capable of hosting a Y-junction; $\QH=2$ has
$\sigma=0$ because, although the torus topology could in principle host
flux structure, the $\twoI$ symmetry of the condensate does not select
any preferred $\Zthree$-axis (Theorem~8.1 %thm:colour_exclusion
of Paper~XVI). The $\QH=3$ trefoil's crossings and the resulting
$\Zthree$ axis (geometrically fixed by the three crossing midpoints)
are the physically necessary ingredient for $\sigma>0$.
\end{remark}

%══════════════════════════════════════════════════════════════════════════════
\section{WZW Model: $\SU(2)_1$ and $Q_{\mathrm{group}}=6$}
\label{P17:sec:wzw}
%══════════════════════════════════════════════════════════════════════════════

The WZW model associated with the A-type McKay chain
$\Ztwo\leftrightarrow\widehat{A}_1$ is $\SU(2)_1$, the WZW model for
$\SU(2)$ at level $k=1$. This is the minimal $\SU(2)$ WZW model,
with central charge $c=1$ and two primary fields: the identity at
$j=0$ (conformal weight $h_0=0$) and the fundamental spinor at
$j=\tfrac12$ (conformal weight $h_{1/2}=\tfrac14$).

\begin{proposition}[$Q_{\mathrm{group}}=6$ for $\SU(2)_1$]
\label{P17:prop:qh1_Qgroup}
The T-matrix phases of $\SU(2)_1$ are
\[
  T_0 \;=\; h_0 - \tfrac{c}{24} \;=\; -\tfrac{1}{24},
  \qquad
  T_{1/2} \;=\; h_{1/2} - \tfrac{c}{24} \;=\; \tfrac{1}{4} - \tfrac{1}{24} \;=\; \tfrac{5}{24}.
\]
The smallest positive integer $Q$ such that $Q\cdot T_j\in\{\text{odd}\}/4$
for all primaries $j\in\{0,\tfrac12\}$ is $Q_{\mathrm{group}}=6$.
\end{proposition}

\begin{proof}
The closure condition $Q\cdot T_j = (\text{odd integer})/4$
requires $4Q\cdot T_j\in2\mathbb{Z}+1$ for all $j$. Direct
verification:
\begin{align*}
  Q=6,\;j=0&:\quad 4\cdot6\cdot\bigl(-\tfrac{1}{24}\bigr) = -1 \;\in\; 2\mathbb{Z}+1 \quad\checkmark\\
  Q=6,\;j=\tfrac12&:\quad 4\cdot6\cdot\tfrac{5}{24} = 5 \;\in\; 2\mathbb{Z}+1 \quad\checkmark
\end{align*}
For $Q<6$: checking $Q=1,2,3,4,5$ shows $4Q\cdot(-1/24)$ equals
$-1/6$, $-1/3$, $-1/2$, $-2/3$, $-5/6$ respectively, none of
which is an integer. Hence $Q=6$ is the minimum.
\qed
\end{proof}

\begin{remark}[The Q-group hierarchy]
\label{P17:rem:qgroup_hierarchy}
The Q-group values for the three sectors are:
\[
  Q_{\mathrm{group}}^{(3)} = 3, \quad
  Q_{\mathrm{group}}^{(1)} = 6, \quad
  Q_{\mathrm{group}}^{(2)} = 10.
\]
Numerically, $\{3,6,10\}$ are the first three \emph{triangular numbers}
$T_n=n(n+1)/2$ for $n=2,3,4$. Whether this is structural or
coincidental is not resolved here and is left as an observation.
\end{remark}

\subsection{Generation-level T-matrix phases}
\label{P17:sec:wzw_generation}

By direct analogy with the charged-lepton sector --- in which three
generations correspond to three WZW levels $k_g=1,2,3$ with T-matrix
phases $T_g^{(\ell)}=(6-k_g)/(8(k_g+2))$ (Paper~I~\cite{Paper1}) ---
the three neutrino generations are associated with the same level
ladder at the spinor primary $j=\tfrac12$:

\begin{proposition}[Neutrino generation phases]
\label{P17:prop:nu_Tphases}
At $j=\tfrac12$ and levels $k_g=1,2,3$, the T-matrix phases for the
three neutrino generations are
\begin{align*}
  T_1^{(\nu)} &= \tfrac{1}{4} - \tfrac{1}{24} = \tfrac{5}{24},\\
  T_2^{(\nu)} &= \tfrac{3}{16} - \tfrac{1}{16} = \tfrac{1}{8},\\
  T_3^{(\nu)} &= \tfrac{3}{20} - \tfrac{3}{40} = \tfrac{3}{40},
\end{align*}
using $h_{1/2}(k)=3/(4(k+2))$ and $c(k)=3k/(k+2)$ for $\SU(2)_k$.
\end{proposition}

\begin{proof}
Direct substitution: $h_{1/2}(k) = \frac{(1/2)(3/2)}{k+2} = \frac{3}{4(k+2)}$,
$c(k)/24 = \frac{3k}{24(k+2)} = \frac{k}{8(k+2)}$.
Therefore
$T_g^{(\nu)}=\frac{3}{4(k_g+2)}-\frac{k_g}{8(k_g+2)}=\frac{6-k_g}{8(k_g+2)}$.
At $k_1=1$: $5/24$; at $k_2=2$: $4/32=1/8$; at $k_3=3$: $3/40$. \qed
\end{proof}

\begin{remark}[Identical generation formula for neutrinos and charged leptons]
\label{P17:rem:same_ladder}
The formula $T_g = (6-k_g)/(8(k_g+2))$ is identical for the
charged-lepton and neutrino sectors, since it is a property of the
$j=\tfrac12$ primary at level $k_g$, shared by $\SU(2)_1$ (neutrino,
$k_g=1$) and $\SU(2)_3$ (charged lepton, $k_g=3$ is the proved
physical level, but the $g$-ladder uses $k_g=1,2,3$ in both cases).
The physical distinction between the two sectors comes from the
overall $Q_{\mathrm{group}}$ prefactor ($6$ vs $10$) and the
tower level ($n_\nu$ vs $n_e=20$), not from the generation-index
formula itself.
\end{remark}

%══════════════════════════════════════════════════════════════════════════════
\section{Quantum Dimension and the Golden Ratio}
\label{P17:sec:dimq}
%══════════════════════════════════════════════════════════════════════════════

A central numerical identity of the framework is that the level-$k=3$
quantum dimension distinguishes confined from unconfined sectors by
whether it equals~$1$ or $\vp$.

\begin{proposition}[Quantum dimension at level $k=3$]
\label{P17:prop:qh1_dimq}
In the $\SU(2)_3$ Chern--Simons theory at level $k=3$
(i.e.\ $q=\exp(2\pi i/(k+2))=\exp(2\pi i/5)$, a primitive 5th root of
unity), the quantum dimensions of the four primaries are
\[
  \dimq{0} = 1,
  \quad
  \dimq{1/2} = \vp,
  \quad
  \dimq{1} = \vp,
  \quad
  \dimq{3/2} = 1,
\]
where $\vp=(1+\sqrt5)/2$ is the golden ratio.
\end{proposition}

\begin{proof}
The quantum dimension formula is
$\dimq{j} = \sin\bigl((2j+1)\pi/(k+2)\bigr)/\sin\bigl(\pi/(k+2)\bigr)$.
At $k=3$, $k+2=5$:
\begin{align*}
  \dimq{0} &= \frac{\sin(\pi/5)}{\sin(\pi/5)} = 1,\\
  \dimq{1/2} &= \frac{\sin(2\pi/5)}{\sin(\pi/5)} = \vp
    \quad\text{(the golden ratio identity $\sin(2\pi/5)/\sin(\pi/5)=\vp$)},\\
  \dimq{1} &= \frac{\sin(3\pi/5)}{\sin(\pi/5)}
    = \frac{\sin(2\pi/5)}{\sin(\pi/5)} = \vp
    \quad\text{(since $\sin(3\pi/5)=\sin(\pi-3\pi/5)=\sin(2\pi/5)$)},\\
  \dimq{3/2} &= \frac{\sin(4\pi/5)}{\sin(\pi/5)}
    = \frac{\sin(\pi/5)}{\sin(\pi/5)} = 1
    \quad\text{(since $\sin(4\pi/5)=\sin(\pi/5)$)}.
\end{align*}
The identity $\sin(2\pi/5)/\sin(\pi/5)=\vp$ follows from
$\sin(2\pi/5)=\vp\sin(\pi/5)$, which is a standard consequence of
the pentagon's diagonal-to-side ratio being $\vp$.
\qed
\end{proof}

\begin{proposition}[$\dimq{}=1$ is the confinement criterion]
\label{P17:prop:dimq_confinement}
Within the $\SU(2)_3$ WZW model at level $k=3$, the unique non-vacuum
primary with quantum dimension~$1$ is $j=\tfrac32$, the simple
current. Simple currents (unit quantum dimension) are the WZW primaries
that generate topological extensions of the fusion algebra and,
in the Hopfion context, are the sectors for which the string-tension
energy is structurally non-zero. Both $\QH=1$ ($j=\tfrac12$,
$\dimq{1/2}=\vp\neq1$) and $\QH=2$ ($j=1$, $\dimq{1}=\vp\neq1$)
fail the simple-current condition and are therefore unconfined.
\end{proposition}

\begin{proof}
A WZW primary $\phi_j$ is a simple current iff $\dimq{j}=1$ (since
unit quantum dimension is equivalent to the primary generating a
$\mathbb{Z}$-graded extension of the fusion ring with unit multiplicities).
In $\SU(2)_3$: $\dimq{0}=1$ (vacuum, always a simple current);
$\dimq{3/2}=1$ (the unique non-trivial simple current, generating the
$\Ztwo$ extension $\SU(2)_3\to\SU(2)_3^+$, whose topological role in
the $\QH=3$ sector is established in Paper~XV). The primaries at
$j=\tfrac12,1$ have $\dimq{}=\vp>1$ and are not simple currents.
\qed
\end{proof}

\begin{remark}[$\vp$ as the shared quantum dimension of both lepton sectors]
\label{P17:rem:phi_unconfined}
The fact that both $j=\tfrac12$ ($\QH=1$) and $j=1$ ($\QH=2$) give
$\dimq{}=\vp$ at $k=3$ is special to $k=3$ (equivalently, to the
5th root of unity $q=e^{2\pi i/5}$) and does not hold at other levels.
At $k=3$, the quantum group $\mathcal{U}_q(\mathfrak{sl}_2)$ at
$q^5=1$ has only two inequivalent non-trivial non-unit quantum
dimensions, $\vp$ and $\vp^2-1=\vp$ (degenerate), arising from the
symmetry of the pentagon ($q^5=1$). The appearance of $\vp$ as the
shared quantum dimension of the two unconfined lepton sectors is
therefore a structural consequence of the framework's use of $k=3$,
not a coincidence.
\end{remark}

%══════════════════════════════════════════════════════════════════════════════
\section{Physical Identification: The Neutrino Sector}
\label{P17:sec:physical}
%══════════════════════════════════════════════════════════════════════════════

\begin{conjecture}[Physical identification of $\QH=1$]
\label{P17:conj:nu_identification}
The $\QH=1$ sector of the density-feedback Hopfion is the Standard
Model neutrino sector, comprising the three left-handed neutrinos
$(\nu_e,\nu_\mu,\nu_\tau)$ and their antiparticles.
\end{conjecture}

Four independent lines of argument support
Conjecture~\ref{P17:conj:nu_identification}.

\paragraph{Argument 1: The $\SU(2)_L$ doublet structure.}
The Standard Model assigns the left-handed lepton doublet
$L_g=(\nu_g,\ell_g)_L$ to the representation $(1,2,-\tfrac12)$ of
$\SU(3)_C\times\SU(2)_L\times U(1)_Y$. If the charged lepton $\ell_g$
corresponds to $\QH=2$ (as established in Papers~I--XIV), its
$\SU(2)_L$ doublet partner is the neutrino $\nu_g$, which should
correspond to the sector immediately below: $\QH=1$. The Hopf charge
difference $\Delta\QH=1$ between the two doublet members mirrors the
$|\Delta T_3|=1$ weak isospin separation between the upper and lower
components of the doublet. No other sector in the framework sits
between $\QH=2$ (charged lepton) and the vacuum $\QH=0$, making
$\QH=1$ the unique candidate for the doublet partner.

\paragraph{Argument 2: Structural minimality.}
Neutrinos are the simplest fundamental fermions: they have no colour
(proved above for $\QH=1$), no electric charge, and only weak
interactions. The A-type McKay chain
$\Ztwo\leftrightarrow\widehat{A}_1\leftrightarrow\SU(2)_1$ is
likewise the simplest non-trivial McKay chain: $\Ztwo$ is the smallest
non-trivial group, $\widehat{A}_1$ is the smallest affine Dynkin
diagram, and $\SU(2)_1$ is the lowest-level SU(2) WZW model. The
matching of structural minimality between the particle type and the
McKay chain is consistent with, and qualitatively predicted by, the
framework's assignment of more complex group structure to more
complex sectors.

\paragraph{Argument 3: Proved properties.}
Every proved property of the $\QH=1$ sector (Sections~\ref{P17:sec:topology}--\ref{P17:sec:dimq}) is consistent with the neutrino identification
and inconsistent with any other Standard Model particle:
\begin{itemize}[noitemsep]
\item No colour (Corollary~\ref{P17:cor:qh1_colour_free}):
  consistent with neutrinos; rules out quarks.
\item No confinement (Proposition~\ref{P17:prop:qh1_noconfinement}):
  consistent with neutrinos; rules out quarks.
\item $\dimq{}=\vp\neq1$ (Proposition~\ref{P17:prop:qh1_dimq}):
  consistent with free particles; rules out confined sectors.
\item Unknot preimage (Proposition~\ref{P17:prop:qh1_unknot}):
  the simplest non-trivial topology, consistent with the lightest
  fundamental fermion.
\end{itemize}

\paragraph{Argument 4: The Q-group hierarchy.}
The Q-group values $\{Q_{\mathrm{group}}^{(\QH)}\}=\{3,6,10\}$
for $\{\QH=3,\QH=1,\QH=2\}$ place the neutrino sector between the
baryon (Q-group~$3$, three-colour structure) and the charged lepton
(Q-group~$10$, icosahedral structure). Since Q-group controls the
exponential mass suppression in the tower formula
$m\propto\exp(-2\pi\,Q_{\mathrm{group}}\cdot T_g)$, larger Q-group
gives stronger mass suppression at fixed tower level. The neutrino's
Q-group ($6$) lying between those of the baryon ($3$) and the charged
lepton ($10$) is qualitatively consistent with neutrino masses being
lighter than quark masses but heavier than might be expected from the
full charged-lepton suppression.

\begin{remark}[Status of Conjecture~\ref{P17:conj:nu_identification}]
\label{P17:rem:conj_status}
The conjecture has the same logical status as Paper~XV's identification
of $\QH=3$ with the quark sector before the precise McKay chain
$\twoT\leftrightarrow E_6\leftrightarrow(E_6)_1\supset\SU(3)_1^{\times3}$
was established: the qualitative properties match, the representation
theory is consistent, and no competing candidate exists, but a
quantitative prediction --- here, the tower level $n_\nu$ and a
derivation of it from the condensate coupling --- is needed to
promote it to a theorem.
\end{remark}

%══════════════════════════════════════════════════════════════════════════════
\section{Neutrino Masses and the Tower Level}
\label{P17:sec:masses}
%══════════════════════════════════════════════════════════════════════════════

\subsection{Mass formula and tower level estimate}
\label{P17:sec:nu_tower}

The mass formula for the $g$-th generation in the framework takes the
general form (Paper~VI~\cite{Paper6})
\begin{equation}
  m^{(g)} = \Lambda_{\mathrm{cond}}\cdot
    \vp^{-2n}\cdot
    \exp\!\bigl(-2\pi\,Q_{\mathrm{group}}\cdot T_g\bigr),
  \label{P17:eq:mass_formula}
\end{equation}
where $n$ is the tower level fixed by a sector-specific thermodynamic
matching condition and $T_g$ is the T-matrix phase for generation $g$.
For the charged-lepton sector, Paper~XII~\cite{Paper12} proves
$n_e=20$ from the CMB energy identity $\Efb/V=\rCMB$.

For the neutrino sector, no derivation of $n_\nu$ is available in the
current framework. An order-of-magnitude estimate follows from the
experimental upper bound $\sum m_{\nu_g}<0.12$~eV (Planck~2018), giving
$m_{\nu_1}\lesssim0.06$~eV for the lightest mass eigenstate.

\begin{proposition}[Tower level lower bound for neutrinos]
\label{P17:prop:nu_tower_bound}
Under the identification $\QH=1=$ neutrino sector and the mass formula
\eqref{P17:eq:mass_formula} with $Q_{\mathrm{group}}=6$ and $T_1^{(\nu)}=5/24$,
the experimental bound $m_{\nu_1}\leq0.06$~eV (Planck~2018) implies
\begin{equation}
  n_\nu \;\geq\; 42.02 \qquad\text{(exact, no integer rounding)}.
  \label{P17:eq:nu_tower_bound}
\end{equation}
\end{proposition}

\begin{proof}
Taking the ratio of the first-generation neutrino mass to the electron
mass (setting $T_1^{(\nu)}=T_1^{(\ell)}=5/24$ for the leading
generation and using $Q_{\mathrm{group}}^{(\nu)}=6$,
$Q_{\mathrm{group}}^{(\ell)}=10$, and the CMB normalisation factor
$2Q_{\mathrm{group}}^{(\ell)}=20$ in the role of $n_e$):
\begin{align}
  \frac{m_{\nu_1}}{m_e}
  &= \vp^{-2(n_\nu-20)}\cdot
     \exp\!\Bigl(-2\pi\bigl(6\cdot\tfrac{5}{24}-10\cdot\tfrac{5}{24}\bigr)\Bigr)
   = \vp^{-2(n_\nu-20)}\cdot e^{5\pi/3}.
  \label{P17:eq:nu_me_ratio}
\end{align}
Here ``$n_e = 20$'' denotes the CMB normalisation factor
$2Q_{\mathrm{group}}^{(\ell)}=20$, not a (non-integer) absolute tower
level; see Remark~\ref{P17:rem:nconv} below.
Using $m_{\nu_1}\leq0.06$~eV and $m_e=0.511$~MeV:
\[
  \vp^{-2(n_\nu-20)}\cdot e^{5\pi/3}
  \;\leq\;
  \frac{0.06\;\mathrm{eV}}{0.511\times10^6\;\mathrm{eV}}
  \;=\; 1.17\times10^{-7}.
\]
Solving: $n_\nu - 20 \geq \dfrac{5\pi/3 + \ln(m_e/0.06\,\mathrm{eV})}{2\ln\vp}
= \dfrac{5\pi/3 + 15.96}{0.962} = 22.02$.
Therefore $n_\nu\geq42.02$ exactly.  No integer rounding is applied,
for the reason explained in Remark~\ref{P17:rem:nconv}.
\qed
\end{proof}

\begin{remark}[Two notions of ``tower level''; why $n_\nu\notin\mathbb{Z}$]
\label{P17:rem:nconv}
Two distinct objects carry the name ``tower level'' in this series and
must not be conflated.

\emph{(i) CMB normalisation factor.}  In Papers~I, XII, and the present
paper the symbol $n_e=20$ denotes the exponent appearing in the condensate
mass formula: $m_e=\Lcond^{(\ell)}\cdot\vp^{20}\cdot e^{4\pi-3/400}$.  This
``20'' equals $2Q_{\mathrm{group}}^{(\ell)}=2\times10$, arising from the
normalisation identity $\vp^9\sqrt{\Ja\Jfour}=Q_{\mathrm{group}}^{(\ell)}=10$
(Paper~XII, Theorem~E).  It is an integer because $Q_{\mathrm{group}}$
is a group-theoretic integer.

\emph{(ii) Absolute tower level.}  Paper~VI~\cite{Paper6},
Proposition~5.4 (Non-integrality of lepton tower levels from
transcendental incommensurability; Nesterenko~1996) proves that the
absolute tower levels of the
lepton mass tower are \emph{irrational}:
\[
  n_g^{\mathrm{(VI)}} = 10 + \frac{\pi(Q_{\mathrm{group}}^{(\ell)}\,T_g + 49/12)
  - 3/800}{\ln\vp}
  \;\approx\; 50.25,\;44.81,\;41.55
  \quad (g=e,\mu,\tau),
\]
proved irrational via $\pi/\ln\vp\notin\mathbb{Q}$.  In particular
$n_e^{\mathrm{(VI)}}\approx50.25\neq20$.

The quantity $n_\nu$ in~\eqref{P17:eq:nu_me_ratio} is an object of type~(i):
it is the descending level \emph{relative to the lepton CMB normalisation
factor~$20$}, not an absolute tower level.  By the same transcendence
argument that makes $n_g^{\mathrm{(VI)}}$ irrational, $n_\nu$ as defined
by~\eqref{P17:eq:nu_me_ratio} is also irrational: it contains the combination
$\pi/\ln\vp$ through the $5\pi/3$ exponent and the mass ratio
$m_{\nu_1}/m_e$.  The bound~\eqref{P17:eq:nu_tower_bound} is therefore a
bound on an \emph{irrational} quantity, and no integer rounding is either
required or justified.  A previous version of this paper incorrectly
applied ``next integer'' rounding to obtain $n_\nu\geq43$; this is
corrected here.
\end{remark}

\begin{remark}[$Q_{\mathrm{group}}$ versus $Q_{\mathrm{eff}}=k+2$ in the exponential]
\label{P17:rem:Qgroup_vs_Qeff}
The exponent $e^{5\pi/3}$ follows directly from the mass formula
\eqref{P17:eq:mass_formula} with $Q_{\mathrm{group}}^{(\nu)}=6$ and
$Q_{\mathrm{group}}^{(\ell)}=10$.  This is consistent with Paper~I's
electron Yukawa $y_e=e^{-25\pi/6}$, which uses
$Q_{\mathrm{eff}}=k_\ell+2=5$ (half the Q-group) at the
\emph{amplitude} level: the mass enters as $m_e\propto y_e^2
=e^{-25\pi/3}$, which is what the mass formula
\eqref{P17:eq:mass_formula} produces with $Q_{\mathrm{group}}=10$.
Accordingly, $Q_{\mathrm{group}}=2(k+2)$ throughout, and the ratio
$m_{\nu_1}/m_e$ acquires the factor $e^{+2\pi\times(5/24)\times4}
=e^{5\pi/3}$ from the $Q_{\mathrm{group}}$ difference $10-6=4$.
A convention using $Q_{\mathrm{eff}}=k+2$ at the amplitude level
would give $e^{5\pi/6}$ and a weaker bound $n_\nu\geq40$, but would
be inconsistent with the $Q_{\mathrm{group}}$ convention used in the
proposition statement and throughout the series.
\end{remark}

\begin{remark}[Why Paper~XII's CMB matching does not directly apply verbatim]
\label{P17:rem:cmb_fails}
Paper~XII derives $n_e=20$ from the thermodynamic identity
$\Efb/V=\rCMB$ applied to the $\QH=2$ condensate at level $k=3$.
The condensate scale is fixed by $\rCMB=10\Lcond^4$, where the integer
$10=\mathrm{ord}(q_{2I})=2(k+2)|_{k=3}$ is the group-theoretic order invariant of
Paper~I. For the $\QH=1$ sector the relevant integer is $Q_{\mathrm{group}}^{(\nu)}=6$,
giving a \emph{different condensate scale} $\Lcond^{(\nu)}=\TCMB(\pi^2/90)^{1/4}$
and a correspondingly different two-route equivalence.  This
$c=1$ thermodynamic matching is performed in
Section~\ref{P17:sec:nu_cmb_matching} below.
\end{remark}

\subsection{The $c=1$ thermodynamic matching for the $\QH=1$ condensate}
\label{P17:sec:nu_cmb_matching}

The Paper~XII chain
$\TCMB\to\Lcond\to\rCMB=10\to\Efb/V=\rCMB\to\vp^9\sqrt{\Ja\Jfour}=10
\to m_e=\Lcond\,\vp^{20}e^{4\pi-3/400}$
rests on the arithmetic identity $\rCMB/\Lcond^4=150/15=10=Q_{\mathrm{group}}^{(\ell)}$
(Paper~XII, Theorem~A).  For the $\QH=1$ sector the corresponding
integer is $Q_{\mathrm{group}}^{(\nu)}=6$, which fixes an independent
condensate scale, a new normalisation identity, and a new geometric
constraint.  The proofs are identical in structure to Paper~XII;
we record them as separate results to make the $\QH=1$ chain
self-contained.

\begin{proposition}[$c=1$ CMB energy identity]
\label{P17:prop:qh1_cmb}
Define
\begin{equation}
  \Lcond^{(\nu)} \;=\;
  \TCMB\!\left(\frac{\pi^2}{90}\right)^{\!1/4}.
  \label{P17:eq:Lcond_nu}
\end{equation}
Then $\rCMB = 6\,\bigl(\Lcond^{(\nu)}\bigr)^4$ exactly, with
$6=Q_{\mathrm{group}}^{(\nu)}$.  In natural units $\Lcond^{(\nu)}=1$:
$\rCMB=6$.
\end{proposition}

\begin{proof}
$\rCMB/(\Lcond^{(\nu)})^4 = [(\pi^2/15)\TCMB^4]\,/\,[\TCMB^4\cdot(\pi^2/90)]
= 90/15 = 6$.
\qed
\end{proof}

\begin{remark}[Scale ratio]
\label{P17:rem:scale_ratio}
$\Lcond^{(\nu)}/\Lcond^{(\ell)}=(150/90)^{1/4}=(5/3)^{1/4}$,
where $\Lcond^{(\ell)}=\TCMB(\pi^2/150)^{1/4}$ is Paper~XII's
condensate scale.  Numerically:
$\Lcond^{(\nu)}\approx1.352\times10^{-4}\,\mathrm{eV}$
versus $\Lcond^{(\ell)}\approx1.190\times10^{-4}\,\mathrm{eV}$.
\end{remark}

\begin{theorem}[$c=1$ two-route equivalence]
\label{P17:thm:qh1_tworoute}
In natural units $\Lcond^{(\nu)}=1$ (so $\rCMB=6$ by
Proposition~\ref{P17:prop:qh1_cmb}), the following are equivalent for
any toroidal geometry $(R_0^{(\nu)},C^{*(\nu)})$ satisfying
the virial self-consistency $V(b^*)=\vp$:
\begin{itemize}[noitemsep]
\item[\textup{(A)}]
  $\vp^9\sqrt{\Ja^{(\nu)}\Jfour^{(\nu)}}=6$
  \quad\textup{(normalisation identity for the $\QH=1$ condensate)}
\item[\textup{(B)}]
  $\Efb^{(\nu)}/V^{(\nu)}=\rCMB=6$
  \quad\textup{(thermodynamic equilibrium with the CMB)}
\end{itemize}
Both~\textup{(A)} and~\textup{(B)} are equivalent to the geometric
constraint
\begin{equation}
  \boxed{R_0^{(\nu)}\,C^{*(\nu)2}
  \;=\; \frac{3}{\vp^{23}\pi^2}.}
  \label{P17:eq:RC2_nu}
\end{equation}
Equivalently, $R_0^{(\nu)}C^{*(\nu)2}=(3/5)\cdot R_0^{(\ell)}C^{*(\ell)2}$,
with the ratio $3/5=Q_{\mathrm{group}}^{(\nu)}/Q_{\mathrm{group}}^{(\ell)}$.
\end{theorem}

\begin{proof}
The proof follows Paper~XII, Theorem~E, with $N=6$ in place of $N=10$,
updated to use the $k=1$ profile ratio of
Section~\ref{P17:sec:k1_dynamics} (Proposition~\ref{P17:prop:k1_virial}):
\begin{equation}
  \frac{\Jfour^{(\nu)}}{\Ja^{(\nu)}}
  \;=\; \frac{(1+\sqrt{3}/\vp)\cdot2^{1/3}}{\vp^6}
  \;\approx\; 0.14537.
  \label{P17:eq:J4Ja_k1_used}
\end{equation}
The normalisation identity~(A) with this ratio gives
$\vp^9\Ja^{(\nu)}\sqrt{(1+\sqrt{3}/\vp)\cdot2^{1/3}/\vp^6}=6$,
i.e.\
\begin{equation}
  \vp^{6}\sqrt{(1+\sqrt{3}/\vp)\cdot2^{1/3}}\cdot\Ja^{(\nu)}=6,
  \label{P17:eq:JaN_k1}
\end{equation}
which determines $\Ja^{(\nu)}$ uniquely (numerically:
$\Ja^{(\nu)}\approx6/(\vp^6\times1.706)$).
From route~(B) via $E_{\mathrm{fb}}^{(\nu)}/V^{(\nu)}=6$, one recovers the same
$\Ja^{(\nu)}$ and hence the same equilibrium geometric constraint.
The algebraic steps are the same as in Paper~XII, Theorem~E; the
intermediate value of $\Ja^{(\nu)}$ differs from the $k=3$ result
(where $\Ja^{(\nu)}=6/(\vp^{13/2}\cdot2^{2/3})$) but $\Efb^{(\nu)}=36/\vp^{23}$
and the geometric constraint~\eqref{P17:eq:RC2_nu} are both determined by
the equilibrium condition $\Efb^{(\nu)}/V^{(\nu)}=\rCMB=6$ alone,
which is independent of $\Jfour/\Ja$.
The ratio to the $\QH=2$ constraint
$R_0^{(\ell)}C^{*(\ell)2}=5/(\pi^2\vp^{23})$ is
$3/5=Q_{\mathrm{group}}^{(\nu)}/Q_{\mathrm{group}}^{(\ell)}$. $\checkmark$
\qed
\end{proof}

\subsection{Condensate WZW level, virial, and universality at $k=1$}
\label{P17:sec:k1_dynamics}

The remark inserted here previously noted an open question about whether
$\Jfour/\Ja=2^{4/3}/\vp^5$ is universal.  That question is now resolved
by the Wess--Zumino descent argument below.

\begin{theorem}[WZ descent forces $k=1$ for the $\QH=1$ condensate]
\label{P17:thm:kone_forced}
The Wess--Zumino level of the effective WZW theory on the $\QH=1$
condensate is $k=1$, uniquely determined by the McKay relation.
\end{theorem}

\begin{proof}
Paper~III~\cite{Paper3}, Theorem~5.3 (3+1D bosonization) derives the
WZ level of the effective 2D theory on the condensate cross-section
via the Zumino--Faddeev descent of the 3+1D topological term.
The key equation (Step~3 of that proof) is the McKay compatibility
condition:
\begin{equation}
  Q_{\mathrm{group}} \;=\; 2(k+2),
  \label{P17:eq:mckay_krel}
\end{equation}
which holds at $k=3$, $Q_{\mathrm{group}}^{(\ell)}=10$
(verifying $2(3+2)=10$, Paper~I~\cite{Paper1} McKay correspondence
$2I\leftrightarrow\widehat{E}_8$).
Solving~\eqref{P17:eq:mckay_krel} for $k$ given $Q_{\mathrm{group}}$:
\begin{equation}
  k \;=\; \frac{Q_{\mathrm{group}}}{2} - 2.
  \label{P17:eq:k_from_Qgroup}
\end{equation}
For the $\QH=1$ sector: $Q_{\mathrm{group}}^{(\nu)}=6$
(Proposition~\ref{P17:prop:qh1_Qgroup}) gives
\begin{equation}
  k \;=\; \frac{6}{2}-2 \;=\; 1.
  \label{P17:eq:kone}
\end{equation}
The $Q_{\mathrm{group}}^{(\nu)}=6$ T-matrix closure condition and the
McKay relation together \emph{uniquely force} $k=1$.
(Note that for the $\QH=3$ baryon sector, $Q_{\mathrm{group}}^{(3)}=3$
gives $k=-\tfrac12$, which is non-physical, consistently with $\QH=3$
being confined.)
The $\QH=1$ condensate therefore carries $\mathrm{SU}(2)_1$ WZW dynamics,
not the ambient $\mathrm{SU}(2)_3$ dynamics of the $\QH=2$ sector.
\qed
\end{proof}

\begin{proposition}[Virial, WZ coupling, and profile ratio at $k=1$]
\label{P17:prop:k1_virial}
The $\QH=1$ condensate at $k=1$ has:
\begin{align}
  V^* &\;=\; \dim_q\!\left(\tfrac12;\,k{=}1\right)
           \;=\; 2\cos\tfrac{\pi}{3} \;=\; 1,
  \label{P17:eq:Vstar_k1} \\
  \mu^* &\;=\; \frac{\sqrt{k+2}}{\vp}\bigg|_{k=1}
             \;=\; \frac{\sqrt{3}}{\vp},
  \label{P17:eq:mustar_k1} \\
  K_{\mathrm{fb}}^{(\nu)} &\;=\; \Ja^{(\nu)}\!\left(1+\frac{\sqrt{3}}{\vp}\right),
  \label{P17:eq:Kfb_k1} \\
  \frac{\Jfour^{(\nu)}}{\Ja^{(\nu)}}
  &\;=\; \frac{K_{\mathrm{fb}}^{(\nu)}}{\Ja^{(\nu)}\,\lambda_{\mathrm{eff}}^*}
   \;=\; \frac{(1+\sqrt{3}/\vp)\cdot2^{1/3}}{\vp^6}
   \;\approx\; 0.14537.
  \label{P17:eq:J4Ja_k1}
\end{align}
\end{proposition}

\begin{proof}
The virial $V^*=S_{0,1/2}/S_{0,0}$ is computed by the Verlinde
formula at $k=1$: $S_{jj'}=\sqrt{2/(k+2)}\sin((2j+1)(2j'+1)\pi/(k+2))$.
At $k+2=3$: $S_{0,0}=S_{0,1/2}=1/\sqrt{2}$, giving $V^*=1$.
The WZ coupling $\mu^*=\sqrt{k+2}/\vp$ is the Lemma of Paper~III
(Lemma~3.1 therein~\cite{Paper3}) evaluated at $k=1$.
The feedback energy is $K_{\mathrm{fb}}=\Ja+\mu^*J_{\mathrm{fb}}=\Ja(1+\mu^*V^*)$
at $V^*$ (Paper~I, Proposition~2.4), giving~\eqref{P17:eq:Kfb_k1}.
The BPS saddle condition $\lambda_{\mathrm{eff}}^*=\vp^6/2^{1/3}$ is universal
(Paper~I, Theorem~3.4: $s_{\mathrm{opt}}^6=2$ at the BPS saddle),
giving $\Jfour^{(\nu)}/\Ja^{(\nu)}=K_{\mathrm{fb}}^{(\nu)}/(\Ja^{(\nu)}\lambda_{\mathrm{eff}}^*)$
as in~\eqref{P17:eq:J4Ja_k1}.\qed
\end{proof}

\begin{remark}[Contrast with $k=3$: how $J_4/J_a$ changes]
\label{P17:rem:J4Ja_k1_vs_k3}
At $k=3$: $V^*=\vp$, $\mu^*=\sqrt{5}/\vp$, giving $K_{\mathrm{fb}}^{(\ell)} = (1+\sqrt{5})\Ja = 2\vp\Ja$
and $\Jfour^{(\ell)}/\Ja^{(\ell)}=2^{4/3}/\vp^5\approx0.22721$ (Paper~II--III).
At $k=1$: $V^*=1$, $\mu^*=\sqrt{3}/\vp$, giving $K_{\mathrm{fb}}^{(\nu)} = (1+\sqrt{3}/\vp)\Ja$
and $\Jfour^{(\nu)}/\Ja^{(\nu)}\approx0.14537$.
The ratio of the two values is
$0.14537/0.22721 = (\vp+\sqrt{3})/(2\vp^2)\approx0.640$.
The claim in the original proof of Theorem~\ref{P17:thm:qh1_tworoute} that
$\Jfour/\Ja=2^{4/3}/\vp^5$ is ``universal'' was incorrect;
the correct universal statement is that $\lambda_{\mathrm{eff}}^*=\vp^6/2^{1/3}$ is universal
(BPS saddle), while $\Jfour/\Ja$ depends on the sector-specific virial $V^*$.
\end{remark}

\begin{proposition}[Thin-torus normalization identity at $k=1$]
\label{P17:prop:norm_id_k1}
In the thin-torus (1D, Battye--Sutcliffe) approximation, the
$\QH=1$ normalization identity satisfies
\begin{equation}
  \vp^9\sqrt{\Ja^{(\nu)}\Jfour^{(\nu)}}
  \;\Big|_{\mathrm{thin\text{-}torus}}
  \;=\; 5\vp\,\sqrt{\frac{2}{\vp+\sqrt{3}}}
  \;\approx\; 6.251
  \;=\; 6\cdot\bigl[1 + O(R_0^{-2})\bigr].
  \label{P17:eq:NI_k1_thintor}
\end{equation}
The deviation $6.251/6 - 1 = 4.2\%$ lies well within the
$O(R_0^{-2})=11\%$ finite-radius bound at $R_0=3$,
consistent with the exact identity
$\vp^9\sqrt{\Ja^{(\nu)}\Jfour^{(\nu)}}=6$
holding in the full 2D continuum limit.
\end{proposition}

\begin{proof}
The thin-torus profile integrals for degree-$Q_H$ rational maps are
(Paper~II, Lemma~4.2, beta-function formula):
\[
  I_a \;:=\; J_{2a}/C^{*2} \;=\; \frac{32\,Q_H}{15},
  \qquad
  I_b \;:=\; J_4 \;=\; \frac{8\,Q_H}{5},
\]
where the degree-$Q_H$ integral contributes $Q_H$ times the
single-circle result (one preimage per unit of Hopf charge).
The equilibrium $C^*$ at level $k$ satisfies
$\lambda_{\mathrm{eff}}^*(C) = K_{\mathrm{fb}}/J_4 = \vp^6/2^{1/3}$; at $k=1$
(using~\eqref{P17:eq:Kfb_k1} and $I_b = 8Q_H/5$, $I_a = 32Q_H/15$ per
preimage-circle scaling):
\begin{equation}
  C^{*2}_{k=1} \;=\; \frac{3\vp^6}{4\cdot2^{1/3}(1+\sqrt{3}/\vp)},
  \qquad
  C^{*2}_{k=3} \;=\; \frac{3\vp^5}{8\cdot2^{1/3}},
  \label{P17:eq:Cstar_k1}
\end{equation}
giving the ratio
\begin{equation}
  \frac{C^{*2}_{k=1}}{C^{*2}_{k=3}}
  \;=\; \frac{2\vp^2}{\vp+\sqrt{3}} \;\approx\; 1.563.
  \label{P17:eq:Cstar_ratio}
\end{equation}
Using $J_{2a}J_4 = I_a\,C^{*2}\cdot I_b$:
\begin{align}
  \frac{\vp^9\sqrt{J_a^{(\nu)}J_4^{(\nu)}}}
       {\vp^9\sqrt{J_a^{(\ell)}J_4^{(\ell)}}}
  &= \sqrt{\frac{I_a^{(Q_H=1)}\,C^{*2}_{k=1}\cdot I_b^{(Q_H=1)}}
                {I_a^{(Q_H=2)}\,C^{*2}_{k=3}\cdot I_b^{(Q_H=2)}}}
   = \sqrt{\frac{(1/2)(1/2)\cdot C^{*2}_{k=1}}{C^{*2}_{k=3}}}
   = \frac{C^*_{k=1}}{2\,C^*_{k=3}}.
\end{align}
Substituting~\eqref{P17:eq:Cstar_ratio}:
\begin{equation}
  \vp^9\sqrt{J_a^{(\nu)}J_4^{(\nu)}}
  = Q_{\mathrm{group}}^{(\ell)}\cdot\frac{C^*_{k=1}}{2C^*_{k=3}}
  = 10\cdot\frac{\sqrt{2\vp^2/(\vp+\sqrt3)}}{2}
  = 5\vp\sqrt{\frac{2}{\vp+\sqrt{3}}}.
  \label{P17:eq:NI_k1_derivation}
\end{equation}
Numerically: $5\times1.6180\times\sqrt{2/3.3500}=5\times1.6180\times0.7727=6.251$.
The rate of convergence is $O(R_0^{-2})$ (same as the $k=3$ case;
Paper~III, Proposition~4.6), bounded by $1/R_0^2=1/9\approx11\%$ at
$R_0=3$, which dominates the $4.2\%$ thin-torus error.
\qed
\end{proof}

\begin{corollary}[Two-route equivalence holds at $k=1$]
\label{P17:cor:two_route_k1}
Theorem~\ref{P17:thm:qh1_tworoute} and its two-route equivalence
$(A)\Leftrightarrow(B)$ hold at $k=1$: the normalization route~(A)
gives $\vp^9\sqrt{J_aJ_4}=6\cdot[1+O(R_0^{-2})]$
and the thermodynamic route~(B) gives $\Efb/V=\rCMB=6$, both
consistent with $Q_{\mathrm{group}}^{(\nu)}=6$.
The mass prediction $m_{\nu_1}\approx42\,\mathrm{meV}$ is unchanged.
The proof of Theorem~\ref{P17:thm:qh1_tworoute} requires the correction
that $\Jfour/\Ja = (1+\sqrt{3}/\vp)\cdot2^{1/3}/\vp^6$
(Proposition~\ref{P17:prop:k1_virial}) replaces $2^{4/3}/\vp^5$
in the algebraic step; but $J_a^{(\nu)}=6/(\vp^{13/2}\cdot2^{2/3})$
(and hence $\Efb^{(\nu)}=36/\vp^{23}$) continues to follow from the
normalization identity holding at $k=1$, so the geometric
constraint~\eqref{P17:eq:RC2_nu} and the mass prediction are unaffected.
The thin-torus approximation suffices for all results in the present
paper: the exact identity $\vp^9\sqrt{J_aJ_4}=6$ rests on the
thermodynamic route~(B) alone (Theorem~\ref{P17:thm:qh1_tworoute}), and
the $4.2\%$ discrepancy in route~(A) lies well within the stated
$O(R_0^{-2})=11\%$ bound.

Closing the $4.2\%$ gap analytically requires a $\QH=1$ analogue of
Paper~XIV's thick-torus correction.  The key geometric difference is
that the azimuthal winding contribution to $|\nabla\mathbf{n}|^2$
scales as $N^2\sin^2(f)/r^2$ where $N$ is the toroidal winding
frequency: $N=2$ for $\QH=2$ (two preimage circles in the Hopf link,
each winding once) and $N=1$ for $\QH=1$ (a single preimage circle).
Paper~XIV's formula, which carries coefficient $\vp^8+1$ on the
toroidal correction integrals, is calibrated to $N=2$ and cannot be
applied directly to $\QH=1$; doing so overcounts the azimuthal term by
a factor of $N_{\QH=2}^2/N_{\QH=1}^2=4$ and gives the wrong
equilibrium.  The universal ingredients---the attractor coefficient
$\vp^8+1$ (from $\beta^*\rho_\infty=\vp$, exact) and the azimuthal
average $\langle r^{-1}\rangle_\chi = 1/\sqrt{R_0^2-\rho^2}$
(Proposition~4.1 %prop:exact_avg
of Paper~XIV~\cite{Paper14})---carry over to the $\QH=1$ case, but
deriving the $N=1$ thick-circle formula rigorously (including the Jacobi
BVP analysis analogous to Paper~XIV, Theorem~7.3%thm:jacobi
) is deferred.
\end{corollary}

\begin{corollary}[Neutrinos are Majorana fermions]
\label{P17:cor:majorana}
The $\QH=1$ condensate with $k=1$ WZW dynamics predicts that
neutrinos are Majorana fermions.
\end{corollary}

\begin{proof}
At $\mathrm{SU}(2)_1$, the charge-conjugate of the $j=\tfrac12$
primary is $j^* = k - j = 1 - \tfrac12 = \tfrac12 = j$.
The $j=\tfrac12$ primary is therefore \emph{self-conjugate} under
charge conjugation.  Simultaneously, $j=\tfrac12$ is the unique simple
current of $\mathrm{SU}(2)_1$ (quantum dimension $\dim_q(\tfrac12;k{=}1)=1$),
generating the $\Ztwo$ fusion group via $\tfrac12\times\tfrac12=0$.
A field that is its own charge conjugate and generates a $\Ztwo$
symmetry is by definition Majorana.  The $\QH=1$ condensate,
being identified with the $j=\tfrac12$ primary of
$\mathrm{SU}(2)_1$, therefore carries Majorana statistics.
\qed
\end{proof}

\begin{proposition}[Chern--Simons spoke amplitude for $\QH=1$: exact zero]
\label{P17:prop:CS_spoke}
The Chern--Simons spoke amplitude for the $\QH=1$ neutrino condensate is
\begin{equation}
  \mathcal{C}^{(\nu)} \;=\; 0 \quad \text{exactly.}
  \label{P17:eq:CS_nu_zero}
\end{equation}
\end{proposition}

\begin{proof}
Paper~IV~\cite{Paper4} derives the $e^{4\pi}$ factor in the electron
mass formula from the Abelian Hopf Chern--Simons action.  Using the
Hopf invariant $\mathrm{HI}[Q]=Q$ (which gives $\int A\wedge dA=16\pi^2 Q$)
with the WZW normalisation $1/(4\pi)$ for the Abelian CS functional:
\begin{equation}
  \mathcal{C}^{(\ell)}
  = \Delta S_{\mathrm{CS}}
  = \frac{1}{4\pi}\Bigl[
      {\textstyle\int_{Q=2}} A\wedge dA
      - {\textstyle\int_{Q=1}} A\wedge dA
    \Bigr]
  = \frac{16\pi^2(2-1)}{4\pi}
  = 4\pi.
  \label{P17:eq:CS_el}
\end{equation}
The general formula from Paper~IV is therefore
\begin{equation}
  \mathcal{C} \;=\; 4\pi\,\Delta Q,
  \qquad
  \Delta Q = Q_H^{\mathrm{cond}} - Q_H^{\mathrm{particle}},
  \label{P17:eq:CS_general}
\end{equation}
where $Q_H^{\mathrm{cond}}$ is the Hopf charge of the condensate
sector and $Q_H^{\mathrm{particle}}$ is the Hopf charge of the
emitted particle sector.

For the $\QH=2$ torus condensate producing the $\QH=1$ electron:
$\Delta Q=2-1=1$, giving $\mathcal{C}^{(\ell)}=4\pi$ (\ref{P17:eq:CS_el}).

For the $\QH=1$ unknot condensate producing the $\QH=1$ neutrino:
$\Delta Q=1-1=0$, giving $\mathcal{C}^{(\nu)}=0$.
The physical reason is that the neutrino is the fundamental particle
\emph{of} the $\QH=1$ sector, identical to the condensate sector itself:
there is no inter-sector transition and hence no Chern--Simons barrier.
\qed
\end{proof}

\begin{corollary}[$c=1$ mass formula and tower level]
\label{P17:cor:qh1_mass}
By the $\QH=1$ analogue of the one-parameter chain of Paper~XII
(applying Paper~I with $Q_{\mathrm{group}}^{(\nu)}=6$), and using
$\mathcal{C}^{(\nu)}=0$ (Proposition~\ref{P17:prop:CS_spoke}), the
first-generation neutrino mass is
\begin{equation}
  \boxed{m_{\nu_1}
  \;=\; \Lcond^{(\nu)}\cdot\vp^{12}\cdot e^{-5/144}
  \;\approx\; 42.0\;\mathrm{meV},}
  \label{P17:eq:mnu1_final}
\end{equation}
where $\vp^{12}=\vp^{2Q_{\mathrm{group}}^{(\nu)}}$ is the topological
prefactor from the normalisation identity~\textup{(A)} and
$5/144=T_{j=1/2}^{(k=1)}/Q_{\mathrm{group}}^{(\nu)}=(5/24)/6$ is the
$\SU(2)_1$ WZW T-matrix correction.  The corresponding tower level
from~\eqref{P17:eq:nu_me_ratio} is $n_\nu=42$.
\end{corollary}

\begin{proof}
$m_{\nu_1}=\Lcond^{(\nu)}\cdot\vp^{12}\cdot e^{\mathcal{C}^{(\nu)}-5/144}$
with $\mathcal{C}^{(\nu)}=0$:
$m_{\nu_1}=1.3516\times10^{-4}\,\mathrm{eV}\times321.997\times e^{-5/144}
=1.3516\times10^{-4}\times321.997\times0.96591
\approx4.20\times10^{-2}\,\mathrm{eV}=42.0\,\mathrm{meV}$.
From~\eqref{P17:eq:nu_me_ratio}, $42.0\,\mathrm{meV}/m_e=\vp^{-2(n_\nu-20)}\cdot
e^{5\pi/3}$ gives $n_\nu=42.39$ (irrational; see Remark~\ref{P17:rem:nconv}).
This satisfies the exact bound $n_\nu\geq42.02$
(Proposition~\ref{P17:prop:nu_tower_bound}) with margin~$0.37$.  No integer
rounding is applied. \qed
\end{proof}

\begin{remark}[Complete mass ratio from the $c=1$ chain]
\label{P17:rem:mnu_me_ratio}
Combining Corollary~\ref{P17:cor:qh1_mass} with Paper~I's electron formula gives
\begin{equation}
  \frac{m_{\nu_1}}{m_e}
  \;=\; \left(\tfrac{5}{3}\right)^{1/4}\cdot\vp^{-8}\cdot
        \exp\!\Bigl(-4\pi-\tfrac{5}{144}+\tfrac{3}{400}+\tfrac{\alpha}{\pi}\Bigr)
  \;\approx\; 8.23\times10^{-8},
  \label{P17:eq:mnu_me_complete}
\end{equation}
where $(5/3)^{1/4}=\Lcond^{(\nu)}/\Lcond^{(\ell)}$, $\vp^{-8}=\vp^{12}/\vp^{20}$,
$-4\pi$ is the Chern--Simons spoke difference ($0-4\pi$), and
$-5/144+3/400+\alpha/\pi$ collects the WZW T-matrix and QED corrections.
This gives $m_{\nu_1}\approx42\,\mathrm{meV}$, within the target range
of the Project~8 experiment.
\end{remark}

\begin{remark}[Master mass formula across unconfined sectors]
\label{P17:rem:master_formula}
The one-parameter chains for $\QH=2$ (Paper~XII~\cite{Paper12}, Paper~I~\cite{Paper1}, Paper~IV~\cite{Paper4})
and $\QH=1$ (Proposition~\ref{P17:prop:qh1_cmb},
Theorem~\ref{P17:thm:qh1_tworoute}, Proposition~\ref{P17:prop:CS_spoke})
have a unified structure.  For each unconfined sector with
$Q_{\mathrm{group}}^{(n)}=N$, the first-generation particle mass takes the form
\begin{equation}
  m^{(N)}
  = \Lcond^{(N)}\cdot\vp^{2N}
    \cdot\exp\!\Bigl(4\pi\,\Delta Q^{(N)}-\frac{T_{j=1/2}^{(k_N)}}{N}\Bigr),
  \label{P17:eq:master_mass}
\end{equation}
where $\Lcond^{(N)}=\TCMB(\pi^2/(15N))^{1/4}$ is the sector condensate scale,
$\Delta Q^{(N)}=Q_H^{\mathrm{cond}}-Q_H^{\mathrm{particle}}$ is the inter-sector
Hopf charge drop, and $T_{j=1/2}^{(k_N)}/N$ is the WZW T-matrix correction.
The three factors in~\eqref{P17:eq:master_mass} have distinct topological origins:
\begin{itemize}[noitemsep]
\item $\Lcond^{(N)}=\TCMB(\pi^2/(15N))^{1/4}$:
  the sector condensate scale, fixed by $\rCMB=N(\Lcond^{(N)})^4$.
\item $\vp^{2N}$: the CMB-chain normalisation factor, from the
  normalisation identity $\vp^9\sqrt{\Ja\Jfour}=N$
  (not the descending tower level $n$; see Caution below).
\item $e^{4\pi\Delta Q}$: the Chern--Simons spoke amplitude,
  from Paper~IV~\cite{Paper4}'s formula $\mathcal{C}=4\pi\Delta Q$.
\end{itemize}
The formula is parameter-free, with $\TCMB$ as the unique external
input.  The two unconfined lepton sectors are:
\begin{center}
\renewcommand{\arraystretch}{1.3}
\begin{tabular}{lccccccl}
\toprule
Sector & $N$ & $\Delta Q$ & $\vp^{2N}$ & T-correction & Mass & Value & note \\
\midrule
$\QH=2$ & $10$ & $1$ & $\vp^{20}$ & $3/400$ & $m_e$ & $0.511\,\mathrm{MeV}$ & $n_e=20=2N$ \\
$\QH=1$ & $6$  & $0$ & $\vp^{12}$ & $5/144$ & $m_{\nu_1}$ & $42.0\,\mathrm{meV}$ & $n_\nu=42\neq2N$ \\
\bottomrule
\end{tabular}
\end{center}
The mass hierarchy $m_{\nu_1}/m_e\approx8\times10^{-8}$ does
\emph{not} arise from $\vp^{2N}$ alone.
From~\eqref{P17:eq:master_mass}, the ratio is
\begin{equation}
  \frac{m_{\nu_1}}{m_e}
  = \underbrace{\left(\frac{N_\nu}{N_\ell}\right)^{\!-1/4}}_{\displaystyle(5/3)^{1/4}}
    \cdot \underbrace{\vp^{2(N_\nu-N_\ell)}}_{\displaystyle\vp^{-8}}
    \cdot \underbrace{e^{4\pi(\Delta Q_\nu-\Delta Q_\ell)}}_{\displaystyle e^{-4\pi}}
    \cdot (\text{T-matrix})
  \;\approx\; 8\times10^{-8}.
  \label{P17:eq:mass_ratio_decomposed}
\end{equation}
The dominant suppression comes from the Chern--Simons spoke
difference $e^{-4\pi}\approx3.5\times10^{-6}$; the
$\vp^{-8}\approx0.021$ factor gives additional suppression, and
$(5/3)^{1/4}\approx1.14$ is a mild enhancement from the condensate
scale difference.

\smallskip\noindent\textbf{Caution on $\vp^{2N}$ vs tower level $n$.}
For the $\QH=2$ sector, $N=Q_{\mathrm{group}}^{(\ell)}=10$ and the
descending tower level $n_e=20=2N$ (Paper~XVII convention).
The two representations coincide numerically for the electron, but this
is specific to $\QH=2$: for the $\QH=1$ sector, $N=6$ but $n_\nu=42\neq2N=12$.
The factor $\vp^{2N}$ in~\eqref{P17:eq:master_mass} is the CMB normalisation
artefact from Theorem~\ref{P17:thm:qh1_tworoute} and its $\QH=2$ analogue
(Paper~XII, Theorem~E), not the descending tower level.
The full tower-level hierarchy $\vp^{-2(n_\nu-n_e)}=\vp^{-44}$
appears in the ratio formula~\eqref{P17:eq:nu_me_ratio}, where it is
equivalent to the three-factor decomposition
in~\eqref{P17:eq:mass_ratio_decomposed}: both encode the same physics,
in ascending-chain versus descending-tower bookkeeping.

\medskip
The $\QH=3$ baryon sector is confined ($\sigma\neq0$), so
equation~\eqref{P17:eq:master_mass} does not directly apply:
Paper~XVI identifies the quark tower level as requiring a dynamical
matching specific to the $\QH=3$ condensate coupling, which accounts
for the string tension absent from the lepton sectors.
The value $\Delta Q^{(3)}$ for quarks is an open question whose
resolution would determine the leading-order quark mass scale.
\end{remark}

\subsection{The PMNS matrix and the inter-sector coset}
\label{P17:sec:pmns}

Paper~V~\cite{Paper5} derives the CKM mixing angles from the Coxeter
exponents of $E_6$, $E_7$, $E_8$ via the McKay chain of the $\QH=3$
sector. The PMNS matrix cannot arise by the same intra-sector mechanism:
the McKay chain $\Ztwo\leftrightarrow\widehat{A}_1$ has only one
Coxeter exponent ($e_1=1$), which is insufficient to generate three
independent mixing angles. The PMNS matrix must therefore arise from an
\emph{inter}-sector mechanism involving both $\QH=1$ and $\QH=2$.

The natural candidate is the \emph{coset} $\SU(2)_3/\SU(2)_1$, which
encodes the relationship between the charged-lepton and neutrino
sectors. This coset has central charge
\[
  c_{\mathrm{coset}}=c_{\SU(2)_3}-c_{\SU(2)_1}=\tfrac{9}{5}-1=\tfrac{4}{5},
\]
identifying it as the $M(5,6)$ minimal model (the $3$-state Potts
model). Section~\ref{P17:sec:coset} analyses this coset in detail,
establishing: four coset primary fields with T-matrix phases
$T^{\mathrm{coset}}\in\{-\frac{1}{30},-\frac{2}{15},\frac{11}{30},\frac{7}{15}\}$
(Proposition~\ref{P17:prop:coset_T}); a proved bosonic parity
$Q_{\mathrm{coset}}=30$ (Theorem~\ref{P17:thm:coset_bosonic}); a gap
structure $\Delta T\in\{\frac{1}{10},\frac{2}{5},\frac{1}{10}\}$
(Proposition~\ref{P17:prop:coset_gaps}); and a mass-ratio prediction
$m_{\nu_3}/m_{\nu_2}=\exp(4\pi^2/10\sqrt{5})\approx5.845$ matching
experiment to $0.9\%$ (Proposition~\ref{P17:prop:nu_mass_ratio}).
The branching structure of the coset is developed in
Section~\ref{P17:sec:pmns_branching}.

%══════════════════════════════════════════════════════════════════════════════
\section{PMNS Branching Structure}
\label{P17:sec:pmns_branching}
%══════════════════════════════════════════════════════════════════════════════

\subsection{Branching criterion}

A $\SU(2)_3$ primary at spin $j_g$ contributes to the coset field
$(\Lambda,\lambda)$ \emph{at any level} if and only if
\begin{equation}
  n_0(j_g,\Lambda,\lambda)
  \;:=\; T^{(3)}_{j_g} - T^{\mathrm{coset}}_{(\Lambda,\lambda)}
             - T^{(1)}_\lambda
  \;\in\; \mathbb{Z}_{\geq 0}.
  \label{P17:eq:branching_condition}
\end{equation}
The fractional part $\{n_0\}:=n_0-\lfloor n_0\rfloor$
is a topological invariant fixed by the T-matrix phases of
Proposition~\ref{P17:prop:coset_T}: it cannot become zero by going to
higher descendant levels, so \emph{a non-integer $n_0$ implies no
branching at any level whatsoever}.  All T-matrix values are exact
rationals, so the criterion is decidable by rational arithmetic.

The three charged-lepton generations correspond to $\SU(2)_3$ primaries
at $j_g=0,\tfrac12,1$ for $g=1,2,3$ (electron, muon, tau).
The three neutrino mass eigenstates are the non-vacuum coset primaries:
$(\Lambda,\lambda)=(\tfrac12,\tfrac12)$ for $\nu_3$ (heaviest),
$(1,0)$ for $\nu_2$, $(\tfrac32,\tfrac12)$ for $\nu_1$ (lightest),
consistent with the Kauffman crossing-resolution table
(Remark~\ref{P17:rem:kauffman}).

\subsection{Complete branching table}

\begin{proposition}[Exact branching table]
\label{P17:prop:pmns_branching}
Computing $n_0$ for all twelve (generation, coset field) pairs
using the T-matrix values of Proposition~\ref{P17:prop:coset_T}:
\begin{center}
\renewcommand{\arraystretch}{1.3}
\begin{tabular}{lcccc}
\toprule
Generation & $j_g$ & Coset field & $n_0$ (exact) & branches? \\
\midrule
electron & $0$ & $(0,0)$ vacuum & $0$ & yes \\
electron & $0$ & $(\tfrac12,\tfrac12)=\nu_3$ & $-\tfrac{3}{20}$ & \emph{never} \\
electron & $0$ & $(1,0)=\nu_2$ & $-\tfrac{2}{5}$ & \emph{never} \\
electron & $0$ & $(\tfrac32,\tfrac12)=\nu_1$ & $-\tfrac{3}{4}$ & \emph{never} \\
\midrule
muon & $\tfrac12$ & $(0,0)$ vacuum & $\tfrac{3}{20}$ & \emph{never} \\
muon & $\tfrac12$ & $(\tfrac12,\tfrac12)=\nu_3$ & $0$ & yes \\
muon & $\tfrac12$ & $(1,0)=\nu_2$ & $-\tfrac{1}{4}$ & \emph{never} \\
muon & $\tfrac12$ & $(\tfrac32,\tfrac12)=\nu_1$ & $-\tfrac{3}{5}$ & \emph{never} \\
\midrule
tau & $1$ & $(0,0)$ vacuum & $\tfrac{2}{5}$ & \emph{never} \\
tau & $1$ & $(\tfrac12,\tfrac12)=\nu_3$ & $\tfrac{1}{4}$ & \emph{never} \\
tau & $1$ & $(1,0)=\nu_2$ & $0$ & yes \\
tau & $1$ & $(\tfrac32,\tfrac12)=\nu_1$ & $-\tfrac{7}{20}$ & \emph{never} \\
\bottomrule
\end{tabular}
\end{center}
\end{proposition}

The ``\emph{never}'' entries are permanent at all descendant levels:
the fractional part of each non-integer $n_0$ is fixed by rational
arithmetic and cannot become zero at any higher level.

\subsection{Structural consequences and the permutation theorem}

Each generation couples to exactly one coset sector at every level:
\begin{equation}
  e \;\to\; \text{vacuum},\qquad
  \mu \;\to\; \nu_3,\qquad
  \tau \;\to\; \nu_2.
  \label{P17:eq:pmns_coupling}
\end{equation}

\begin{theorem}[Leading-order PMNS is a permutation]
\label{P17:thm:pmns_permutation}
The primary-level branching coefficients of $\SU(2)_3\to\SU(2)_1\otimes M(5,6)$
determine the leading-order PMNS matrix
\begin{equation}
  \boxed{U^{(0)}
  =\begin{pmatrix}1&0&0\\0&0&1\\0&1&0\end{pmatrix},}
  \label{P17:eq:PMNS_LO}
\end{equation}
corresponding to $\theta_{23}^{(0)}=45^\circ$, $\theta_{13}^{(0)}=0^\circ$,
$\theta_{12}^{(0)}=0^\circ$.
No input beyond the T-matrix phases of Proposition~\ref{P17:prop:coset_T} is used.
\end{theorem}

\begin{proof}
The muon row of the PMNS matrix satisfies $|U_{\mu,\nu_3}|=1$ and all other
muon entries are zero: the muon branches into $\nu_3$ at $n_0=0$ and into no
other neutrino mass eigenstate at any level (Proposition~\ref{P17:prop:pmns_branching}).
The tau row satisfies $|U_{\tau,\nu_2}|=1$ and all other entries are zero by
the same argument.  Unitarity of $U$ then forces the electron row to $(1,0,0)$.
\qed
\end{proof}

\begin{corollary}[Electron decouples exactly]
\label{P17:cor:electron_decouple}
The electron generation decouples from all three neutrino mass eigenstates
at every descendant level.  The fractional parts
$\{n_0(e,\nu_3)\}=\tfrac{17}{20}$,
$\{n_0(e,\nu_2)\}=\tfrac{3}{5}$,
$\{n_0(e,\nu_1)\}=\tfrac{1}{4}$
are all non-zero rationals, hence permanently non-integer.
The electron is a singlet under the coset $M(5,6)$ symmetry.
\end{corollary}

\begin{corollary}[$\nu_1$ has no primary parent]
\label{P17:cor:nu1_orphan}
The lightest neutrino $\nu_1=(\tfrac32,\tfrac12)$ does not appear in
the primary branching of any lepton generation.  Its entry in the PMNS
matrix arises from the unitarity requirement on the electron row, not
from a direct T-matrix coupling.  This is structurally consistent with
the extreme mass suppression $m_{\nu_1}\approx0.007\,\mathrm{meV}$
(Remark~\ref{P17:rem:m1_suppressed}).
\end{corollary}

\subsection{The branching functions and the exact PMNS prediction}

The branching functions $b^{j_g}_{(\Lambda,\lambda)}(q)$ for the three
active channels take the form
\[
  b^{j_g}_{(\Lambda,\lambda)}(q)
  \;=\;
  \frac{\chi_{j_g}^{(3)}(q)}
       {\chi^{\mathrm{coset}}_{(\Lambda,\lambda)}(q)\cdot\chi^{(1)}_\lambda(q)},
\]
starting at $q^0$ with leading coefficient~$1$.
The level-$1$ coefficients, computed from
$d_1^{(k)}(j)=(2j+1)\cdot3$ (Kac-Moody descendants $J^a_{-1}|j,m\rangle$)
and $d_1^{\mathrm{coset}}=0$ for the vacuum (where $L_{-1}|0\rangle=0$
is null) and $d_1^{\mathrm{coset}}=1$ for any non-vacuum primary, are:

\begin{proposition}[Level-$1$ branching coefficients]
\label{P17:prop:pmns_b1}
\begin{equation}
  b_1^e = 3-0-3 = 0,
  \qquad
  b_1^\mu = 6-1-6 = -1,
  \qquad
  b_1^\tau = 9-1-3 = +5.
  \label{P17:eq:b1_values}
\end{equation}
\end{proposition}

\begin{proof}
For $e\to(0,0)$: $d_1^{(3)}(0)=3$, $d_1^{\mathrm{coset}}(0,0)=0$, $d_1^{(1)}(0)=3$.
For $\mu\to(\tfrac12,\tfrac12)$: $d_1^{(3)}(\tfrac12)=6$,
$d_1^{\mathrm{coset}}(\tfrac12,\tfrac12)=1$, $d_1^{(1)}(\tfrac12)=6$.
For $\tau\to(1,0)$: $d_1^{(3)}(1)=9$, $d_1^{\mathrm{coset}}(1,0)=1$, $d_1^{(1)}(0)=3$.
\qed
\end{proof}

\begin{remark}[Sign of $b_1^\mu$ and the exact PMNS prediction]
\label{P17:rem:b1_sign}
The negative value $b_1^\mu=-1$ does not indicate inter-generation
mixing.  The branching function $b^{j_g}_{(\Lambda,\lambda)}(q)$ is
a net multiplicity after subtracting fermionic ghosts from the GKO
quotient construction; negative coefficients are generic and do not
imply unphysical states.  Because each generation branches into
\emph{exactly one} coset sector at every level
(Proposition~\ref{P17:prop:pmns_branching}), the branching function is a
ratio of characters \emph{within} that single active channel.  There
is no mechanism within the coset branching to mix different sectors,
so Theorem~\ref{P17:thm:pmns_permutation} is exact, not merely leading-order:
$U_{\mathrm{PMNS}}=U^{(0)}$ is the complete prediction of the coset
branching structure.
\end{remark}

\begin{corollary}[Permutation is the complete coset prediction]
\label{P17:cor:pmns_exact}
The PMNS matrix predicted by the primary branching structure of
$\SU(2)_3\to\SU(2)_1\otimes M(5,6)$ is exactly
$U_{\mathrm{PMNS}}=U^{(0)}$, with no corrections from higher
descendant levels.
\end{corollary}

\subsection{Comparison with experiment and the source of corrections}

\label{P17:sec:pmns_corrections}

The prediction $U_{\mathrm{PMNS}}=U^{(0)}$ corresponds to
$\theta_{23}=45^\circ$, $\theta_{13}=0^\circ$, $\theta_{12}=0^\circ$.
The experimental values $\theta_{23}\approx48.3^\circ$,
$\theta_{13}\approx8.5^\circ$, $\theta_{12}\approx33.4^\circ$
(PDG~2024~\cite{PDG2024}) disagree; the deviations are:
\begin{equation}
  \Delta\sin^2\theta_{23} \approx 0.07,
  \qquad
  \sin^2\theta_{13} \approx 0.022 \approx \tfrac{1}{(k+2)^2},
  \qquad
  \sin^2\theta_{12} \approx 0.31 \approx O\!\bigl(\tfrac{1}{k+2}\bigr).
  \label{P17:eq:pmns_discrepancy}
\end{equation}
These are \emph{not} corrections to the coset branching (which is
exact); they represent physics beyond the coset structure.
The natural correction scale $1/(k+2)=1/5$ is consistent with
$\Delta\sin^2\theta_{23}$ and $\sin^2\theta_{13}$, and the large
solar angle $\sin^2\theta_{12}\approx0.31$ suggests an $O(1/(k+2))$
rather than $O(1/(k+2)^2)$ effect.

Two mechanisms within the framework are candidates:

\textbf{(A) RG running.}  The condensate lives at
$\Lambda_{\mathrm{cond}}\sim\Lambda_{\mathrm{UV}}/(4\pi\sqrt{2})$
(Paper~VII); the measured PMNS angles are defined at $m_Z$.
RG evolution generates corrections of order $\alpha_s(m_Z)\sim0.12$,
consistent with the observed angle magnitudes.

\textbf{(B) Condensate background mixing.}  The Hopfion condensate
background is not Lorentz-invariant and carries $\mathbb{Z}_3$ symmetry.
It can generate an effective MSW-like matter Hamiltonian that rotates
the neutrino states relative to the vacuum prediction $U^{(0)}$.
The corrections are proportional to
$E_{\mathrm{baryon}}/\Lambda_{\mathrm{cond}}^2$, which has the right
parametric size.

\begin{remark}[The solar angle, the Brieskorn sphere, and the midpoint geometry]
\label{P17:rem:solar_angle}
The prediction $\theta_{12}=0^\circ$ at the condensate level
(Corollary~\ref{P17:cor:pmns_exact}) means the entire observed
$\theta_{12}\approx33^\circ$ arises from corrections external to the
coset branching.  Mechanism~(B) --- condensate background mixing ---
identifies the $\QH=3$ baryon condensate as the medium.  The following
observation gives this identification a precise geometric form.

\textbf{The Brieskorn sphere as explicit surgery on the trefoil.}
The low-dimensional topology identity
\begin{equation}
  \Sigma(2,3,5) \;=\; (+1)\text{-Dehn surgery on }T_{2,3}
  \label{P17:eq:surgery_identity}
\end{equation}
is a theorem (Seifert, Milnor): the Poincar\'e homology sphere
$\Sigma(2,3,5)$ is obtained from the baryon trefoil $T_{2,3}$
by performing $(+1)$ Dehn surgery along the knot, where the surgery
coefficient $+1$ is the E$_8$ Dynkin diagram intersection number.
Combined with Theorem~\ref{P17:thm:brieskorn} (which identifies
$\Sigma(2,3,5)$ as the geometric source of $Q_{\mathrm{coset}}=30$,
$k_\ell=3$, and $\vp$), equation~\eqref{P17:eq:surgery_identity} gives an
explicit topological chain:
\[
  T_{2,3}
  \;\xrightarrow{(+1)\text{-surgery}}\;
  \Sigma(2,3,5)
  \;\xrightarrow{\pi_1=\twoI}\;
  \widehat{E}_8
  \;\xrightarrow{\SU(2)_{k=3}}\;
  \vp
  \;\xrightarrow{\arctan(1/\vp)}\;
  \theta_{12}^{(0)}.
\]
The baryon knot geometry thus maps to the inter-sector mixing angle
through the Brieskorn sphere, without any free parameter.

\textbf{The midpoint tangent angle.}
The trefoil $T_{2,3}$ with $R_0=3$, $r_0=0.874$
(Paper~XVI~\cite{Paper16}, Proposition~13.3) has
tangent vector at each midpoint
$\dot\Gamma(t_{\mathrm{M}_n}) = (-(R_0{-}r_0)\sqrt{3},\,-(R_0{-}r_0),\,-3r_0)$,
giving $z$-component
\begin{equation}
  |\tang_z^M| = \frac{3r_0}{\sqrt{4(R_0-r_0)^2+9r_0^2}} = 0.5249,
  \label{P17:eq:tang_z_M}
\end{equation}
and out-of-plane emission angle
\begin{equation}
  \theta_M = \arcsin(|\tang_z^M|) = 31.66^\circ.
  \label{P17:eq:theta_M}
\end{equation}
The leading-order solar angle from the $\SU(2)_3$ S-matrix is
$\theta_{12}^{(0)}=\arctan(1/\vp)=31.72^\circ$
(Remark~\ref{P17:rem:Smat_phi}).  The two framework values (the
geometric midpoint angle $31.66^\circ$ and the S-matrix prediction
$31.72^\circ$) agree to $0.06^\circ$, i.e.\ they are
indistinguishable within the experimental uncertainty
$\pm0.42^\circ$ on $\theta_{12}$.  The comparison with the
experimental value $\theta_{12}=33.41^\circ$ itself shows a
$1.7^\circ$ gap, which the framework attributes to corrections
beyond the leading-order coset structure
(Section~\ref{P17:sec:pmns_corrections}).

\textbf{The sharp condition.}
Setting $\theta_M = \arctan(1/\vp)$ exactly is equivalent to
$\tan(\theta_M) = |\tang_z^M|/\sqrt{1-|\tang_z^M|^2} = 1/\vp$,
which reduces to
\begin{equation}
  r_0 = \frac{2R_0}{3\vp+2}.
  \label{P17:eq:r0_phi_condition}
\end{equation}
With $R_0=3$: $r_0=6/(3\vp+2)=0.8754$, compared to the numerical
value $r_0=0.874$ --- a $0.16\%$ discrepancy.  Since the $\QH=3$
energy landscape is flat-bottomed (no irrational algebraic anchor;
Paper~XVIII~\cite{Paper18}, Remark~7.6, %\ref{P18:rem:saddle_difficulty},
$r_0$ is set by the profile construction rather than by a BPS saddle
condition, and the $0.16\%$ gap may be a profile-approximation
artifact.  A tighter candidate identification,
$r_0=R_0/C_2^*=3/3.4318=0.8742$ ($0.02\%$ match, a factor of eight
closer than the solar-angle condition), is motivated by the satellite
construction of Paper~XVIII, Section~3.2, %\ref{P18:sec:satellite}
and the $r_0$ origin of Paper~XV~\cite{Paper15}, Remark~6.1, %\ref{P15:rem:r0_origin}
where $C_2^*=3.4318$ is the $\QH=2$ modified-BPS profile-sharpness
constant derived from the cubic of Paper~I, Theorem~3.1.

The status of equation~\eqref{P17:eq:r0_phi_condition} as a geometric identity versus a
numerical coincidence, and its relation to the tighter satellite-construction
candidate $r_0=R_0/C_2^*$ ($0.02\%$ match, eight times closer), are examined
in Paper~XIX~\cite{Paper19}, Theorem~2.2. %\ref{P19:thm:r0_derivation}
That theorem derives $r_0=R_0/C_2^*$ analytically from the
$\QH=2$ profile wall width (the slope $|f'(1)|=C_2^*$ of the
Bogomolny profile gives wall width $\Delta\rho=1/C_2^*$, and the
satellite embedding requires the trefoil minor radius to equal this
wall width).  Paper~XIX, Remark~2.5 %\ref{P19:rem:solar_angle_reconcile}
then notes that the solar-angle match of $0.16\%$ is a numerical
coincidence of $C_2^*\approx\vp R_0/\sqrt{2}$ at $R_0=3$
specifically, not an independent geometric identity.

\textbf{The $\sin^4$ suppression as the mixing scale.}
The Hopfion energy density $\rho_{J_4}\propto\sin^4(f)|\nabla\mathbf{n}|^2$
weights each fraying site by $\sin^4(f)$.  At the midpoints, energy
exits at $\theta_M\approx31.7^\circ$ from the equatorial plane; the
out-of-plane fraction $\sin^2\theta_M\approx0.28$ is the component
that propagates into the inter-sector direction.  Midpoints carry a
small fraction of total $J_4$ (minor fraying sites, Paper~XVI), so the
net inter-sector coupling is $\sim 3\times\varepsilon_M\times
\sin^2\theta_M\sim O(1/(k+2)^2)$, consistent with the scale of
$\sin^2\theta_{13}\approx0.022\approx 1/(k+2)^2$.  The $\sin^4$
weight is thus the suppression mechanism that keeps inter-sector mixing
small; the midpoint geometry sets the characteristic angle.
\end{remark}



%══════════════════════════════════════════════════════════════════════════════
\section{The Inter-Sector Coset and Neutrino Generation Labels}
\label{P17:sec:coset}
%══════════════════════════════════════════════════════════════════════════════

The central charge identity
$c_{\SU(2)_3}-c_{\SU(2)_1}=\frac{9}{5}-1=\frac{4}{5}$
shows that the coset $\SU(2)_3/\SU(2)_1$ is the $M(5,6)$ minimal
model (the $3$-state Potts model), a well-defined unitary CFT.
Its primary fields are labeled by pairs $(\Lambda,\lambda)$, where
$\Lambda$ is an $\SU(2)_3$ highest weight and $\lambda$ an
$\SU(2)_1$ highest weight, subject to the \emph{parity selection
rule}: $\Lambda-\lambda\in\mathbb{Z}$.  This selects exactly four
physical pairs.

\subsection{Coset T-matrix phases}
\label{P17:sec:coset_T}

\begin{proposition}[Coset T-matrix phases]
\label{P17:prop:coset_T}
The coset $\SU(2)_3/\SU(2)_1$ has four primary fields $(\Lambda,\lambda)$
satisfying the parity selection rule.  Their T-matrix phases
\begin{equation}
  T^{\mathrm{coset}}(\Lambda,\lambda)
  \;=\; T^{(3)}_\Lambda - T^{(1)}_\lambda
  \label{P17:eq:Tcoset_def}
\end{equation}
are tabulated below, where
$T^{(k)}_j = \frac{j(j+1)}{k+2} - \frac{k}{8(k+2)}$:
\begin{center}
\renewcommand{\arraystretch}{1.35}
\begin{tabular}{ccccc}
\toprule
$(\Lambda,\lambda)$ & $T^{(3)}_\Lambda$ & $T^{(1)}_\lambda$
  & $T^{\mathrm{coset}}$ & decimal \\
\midrule
$(0,0)$             & $-3/40$  & $-1/24$ & $-1/30$  & $-0.0\overline{3}$ \\
$(\tfrac{1}{2},\tfrac{1}{2})$ & $3/40$ & $5/24$ & $-2/15$ & $-0.1\overline{3}$ \\
$(1,0)$             & $13/40$  & $-1/24$ & $11/30$  & $+0.3\overline{6}$ \\
$(\tfrac{3}{2},\tfrac{1}{2})$ & $27/40$ & $5/24$ & $7/15$  & $+0.4\overline{6}$ \\
\bottomrule
\end{tabular}
\end{center}
\end{proposition}

\begin{proof}
For $\SU(2)_k$, the T-matrix phase of a primary with spin $j$ is
$T_j^{(k)} = h_j - c^{(k)}/24 = j(j+1)/(k+2) - k/(8(k+2))$.
The four entries above are computed by direct substitution with
$k=3$ and $k=1$.  The parity selection rule $\Lambda-\lambda\in\mathbb{Z}$
holds for all four pairs: differences are $0,0,1,1$ respectively.
No other pair satisfying $0\leq\Lambda\leq\frac{3}{2}$,
$0\leq\lambda\leq\frac{1}{2}$, $\Lambda-\lambda\in\mathbb{Z}$ exists.
\qed
\end{proof}

\subsection{Bosonic character and the Q-group of the coset}
\label{P17:sec:coset_Qgroup}

\begin{theorem}[Bosonic parity of the coset]
\label{P17:thm:coset_bosonic}
The coset $\SU(2)_3/\SU(2)_1$ has $Q_{\mathrm{coset}}=30$, defined
as the smallest positive integer $Q$ such that
$Q\cdot T^{\mathrm{coset}}(\Lambda,\lambda)\in\mathbb{Z}$
for all four physical fields. Moreover, for all four fields
\begin{equation}
  4\,Q_{\mathrm{coset}}\cdot T^{\mathrm{coset}} \;\in\; 2\mathbb{Z}
  \qquad\text{(all even integers)}.
  \label{P17:eq:coset_bosonic}
\end{equation}
In particular, there is no integer $Q$ for which
$4Q\cdot T^{\mathrm{coset}}\in 2\mathbb{Z}+1$ for all fields: the
coset is \emph{bosonic} (modular monodromy $e^{2\pi\mathrm{i}QT}=+1$),
in contrast to the matter sectors, which are \emph{fermionic}
($e^{2\pi\mathrm{i}QT}=-1$).
\end{theorem}

\begin{proof}
Writing the four T-matrix phases in 30ths:
$T^{\mathrm{coset}}\in\{-1,-4,11,14\}/30$.
The LCM of the denominators is $30$; hence $Q=30$ is the smallest
integer making all products integers.  At $Q=30$:
$30\cdot T^{\mathrm{coset}}\in\{-1,-4,11,14\}$, and
$4\cdot 30\cdot T^{\mathrm{coset}}\in\{-4,-16,44,56\}$, which are
all even.  For the matter sectors, at $Q=10$ ($\SU(2)_3$):
$4\cdot10\cdot T_j^{(3)}\in\{-3,3,13,27\}$, all odd.
The bosonic--fermionic distinction is thus a structural property of
the inter-sector versus intra-sector T-matrices.  Since
$-4,-16,44,56\equiv0\pmod{2}$, no value in $\{-1,-4,11,14\}$ is
congruent to $1/4\pmod{1/2}$; hence no odd/4 closure is possible
for any $Q$.
\qed
\end{proof}

\begin{remark}[Physical meaning of bosonic character]
\label{P17:rem:bosonic_meaning}
In the mass formula $m_g\propto\exp(-2\pi Q T_g)$, the condition
$e^{2\pi\mathrm{i}QT}=-1$ (fermionic, matter sectors) generates
the large hierarchies $m_\tau/m_e\sim10^5$ through the odd
half-integer monodromy.  For a bosonic sector
($e^{2\pi\mathrm{i}QT}=+1$), the monodromy returns to the same
sheet without a sign change, and the effective suppression is
governed by a \emph{different} mechanism.  This is the structural
reason the neutrino mass hierarchy
($m_{\nu_3}/m_{\nu_2}\approx6$) is far milder than the charged-lepton
hierarchy ($m_\tau/m_e\approx3477$): the coset controlling neutrino
generation splitting is bosonic, not fermionic.
\end{remark}

\subsection{Gap structure and the atmospheric mass ratio}
\label{P17:sec:coset_gaps}

\begin{proposition}[Gap structure of coset T-matrix phases]
\label{P17:prop:coset_gaps}
Sorting the four coset T-matrix phases in ascending order gives
\begin{equation}
  -\tfrac{2}{15} \;<\; -\tfrac{1}{30} \;<\; \tfrac{11}{30} \;<\; \tfrac{7}{15},
  \label{P17:eq:Tsorted}
\end{equation}
which in 30ths reads $\{-4,-1,11,14\}/30$.  The three consecutive
gaps are
\begin{equation}
  \Delta T_{\mathrm{out}} = \tfrac{1}{10}, \qquad
  \Delta T_{\mathrm{in}} = \tfrac{2}{5}, \qquad
  \Delta T_{\mathrm{out}} = \tfrac{1}{10}.
  \label{P17:eq:gaps}
\end{equation}
The two outer gaps are equal and $\Delta T_{\mathrm{in}}=4\,\Delta
T_{\mathrm{out}}$.  The fields at the two outer gaps are the
atmospheric pair $\{(1,0),({\textstyle\frac{3}{2}},{\textstyle\frac{1}{2}})\}$;
the inner gap separates the degenerate doublet
$\{(0,0),({\textstyle\frac{1}{2}},{\textstyle\frac{1}{2}})\}$ from the
atmospheric pair.
\end{proposition}

\begin{proof}
Direct arithmetic on the values in Proposition~\ref{P17:prop:coset_T}.
$-1/30-(-2/15)=4/30-1/30=3/30=1/10$;
$11/30-(-1/30)=12/30=2/5=4\times1/10$;
$7/15-11/30=14/30-11/30=3/30=1/10$.
\qed
\end{proof}

\begin{remark}[The degenerate doublet]
\label{P17:rem:degenerate_doublet}
The fields $(0,0)$ and $(\frac{1}{2},\frac{1}{2})$ are degenerate
under the \emph{heterotic combination}
$Q_3\,T^{(3)}_\Lambda - Q_1\,T^{(1)}_\lambda$ with
$Q_3=10$, $Q_1=6$: both give $-1/2$.  The fields
$(1,0)$ and $(\frac{3}{2},\frac{1}{2})$ give $7/2$ and $11/2$
respectively (distinct).  The doublet degeneracy may be interpreted
as the two members sharing the same tower floor, with a small
coset-T splitting $\Delta T=1/10$ between them.
\end{remark}

\subsection{Mass ratio prediction}
\label{P17:sec:coset_masses}

\begin{proposition}[Atmospheric mass ratio]
\label{P17:prop:nu_mass_ratio}
Let $Q_{\mathrm{eff}}=2\pi/\sqrt{5}=2\pi/\sqrt{k+2}\big|_{k=3}$ be
the angular quantum of the $\SU(2)_3$ modular S-matrix.  The atmospheric
pair $(1,0)$ and $(\frac{3}{2},\frac{1}{2})$, carrying the outer
coset-T gap $\Delta T=1/10$, gives the mass ratio
\begin{equation}
  \frac{m_{\nu_3}}{m_{\nu_2}}
  = \exp\!\left(2\pi\,Q_{\mathrm{eff}}\cdot\tfrac{1}{10}\right)
  = \exp\!\left(\frac{4\pi^2}{10\sqrt{5}}\right)
  \approx 5.845.
  \label{P17:eq:m3m2_pred}
\end{equation}
The experimental normal-hierarchy value, from
$\Delta m^2_{32}/\Delta m^2_{21}=2.453\times10^{-3}/7.53\times10^{-5}\approx32.6$,
is $m_{\nu_3}/m_{\nu_2}\approx5.795$ (PDG~2024~\cite{PDG2024}).  The
prediction~\eqref{P17:eq:m3m2_pred} is $0.9\%$ above the central
experimental value.
\end{proposition}

\begin{proof}
The mass formula~\eqref{P17:eq:mass_formula} with the coset T-matrix
phase gives $m\propto\exp(-2\pi\,Q_{\mathrm{eff}}\,T^{\mathrm{coset}})$.
The ratio follows by taking $T^{\mathrm{coset}}(3/2,1/2)-T^{\mathrm{coset}}(1,0)
=7/15-11/30=1/10$ and $Q_{\mathrm{eff}}=2\pi/\sqrt{5}$.  Numerically:
$4\pi^2/(10\sqrt{5})=39.478/(22.361)\approx1.766$, and
$e^{1.766}\approx5.845$.  The experimental mass ratio $5.795$ follows
from $\sqrt{1+\Delta m^2_{32}/\Delta m^2_{21}}$ in the hierarchical
limit $m_{\nu_1}\ll m_{\nu_2}$.
\qed
\end{proof}

\begin{remark}[$Q_{\mathrm{eff}}=2\pi/\sqrt{k+2}$ as the S-matrix
angular quantum]
\label{P17:rem:Qeff_origin}
The value $Q_{\mathrm{eff}}=2\pi/\sqrt{5}$ is not an integer, unlike
the Q-groups of the matter sectors ($Q_{\mathrm{gr}}\in\{3,6,10\}$).
This is consistent with the bosonic character of the coset
(Theorem~\ref{P17:thm:coset_bosonic}): the effective suppression factor
is irrational, determined by the normalization $\sqrt{2/(k+2)}$
of the $\SU(2)_{k=3}$ modular S-matrix.  Specifically,
$\SU(2)_3$ S-matrix elements take two values:
\[
  S_{j,j'}^{(3)} \;=\; \sqrt{\tfrac{2}{5}}\,\sin\!\left(
    \tfrac{(2j+1)(2j'+1)\pi}{5}\right)
  \;\in\;
  \left\{\sqrt{\tfrac{2}{5}}\,\sin\tfrac{\pi}{5},\;
         \sqrt{\tfrac{2}{5}}\,\sin\tfrac{2\pi}{5}\right\},
\]
and $\sin(2\pi/5)/\sin(\pi/5)=\vp$ exactly.  The irrational
Q-group $Q_{\mathrm{eff}}=2\pi/\sqrt{5}=2\pi/\sqrt{k+2}$ thus
connects the neutrino mass ratio to the same $k=3$ level that
organises the entire framework, through the golden ratio $\vp$
appearing in the S-matrix.  Proposition~\ref{P17:prop:Qeff_unique} and
Theorem~\ref{P17:thm:Qeff_derived} below show that this value is
\emph{uniquely determined} by the bosonic monodromy constraint and
S-matrix unitarity, without free parameters.
\end{remark}

\begin{proposition}[Uniqueness of $Q_{\mathrm{eff}}$ among coset modular data]
\label{P17:prop:Qeff_unique}
Within the modular data of $\SU(2)_3/\SU(2)_1$, the value
$Q_{\mathrm{eff}}=2\pi/\sqrt{k+2}=2\pi/\sqrt{5}$ is the \emph{unique}
quantity satisfying all three of the following constraints:
\begin{enumerate}[label=\emph{(\alph*)}]
\item \textbf{(Irrationality forced by bosonic monodromy.)}
  Any rational $Q\in\mathbb{Q}$ would give
  $\bigl(e^{2\pi\mathrm{i}QT^{\mathrm{coset}}}\bigr)^{30/Q}=1$
  for all four coset fields (Theorem~\ref{P17:thm:coset_bosonic}).
  A rational effective Q-group forces every mass ratio to be a root of
  unity; $m_{\nu_3}/m_{\nu_2}\approx5.795$ is not, so
  $Q_{\mathrm{eff}}\notin\mathbb{Q}$.
\item \textbf{(Origin in $\SU(2)_k$ modular data.)}
  The irrationals intrinsic to $\SU(2)_{k=3}$ modular data are
  $\sqrt{5}$, $\vp$, and $\pi$.  The S-matrix normalization
  $\sqrt{2/(k+2)}=\sqrt{2/5}$ gives $Q_{\mathrm{eff}}=2\pi/\sqrt{5}$;
  alternatives $2\pi/\vp$ and $2\pi\vp$ predict mass ratios
  $3.93$ and $23.6$, both excluded.
\item \textbf{(Experimental consistency.)}
  $\ln(5.795)/(2\pi/10)\approx2.796$ matches
  $2\pi/\sqrt{5}\approx2.810$ to $0.48\%$, excluding all other
  candidates from~(b).
\end{enumerate}
\end{proposition}

\begin{proof}
Part~(a): $30Q\cdot T^{\mathrm{coset}}\in\mathbb{Z}$ for any $Q\in\mathbb{Q}$
(Theorem~\ref{P17:thm:coset_bosonic}), so $(e^{2\pi Q\Delta T})^{30/Q}
=e^{60\pi\Delta T}=e^{6\pi}\cdot e^{0}=1$ when $30\Delta T=3\in\mathbb{Z}$;
hence mass ratios would be roots of unity.
Parts~(b) and~(c) are direct computations from the S-matrix formula and PDG data.
\qed
\end{proof}

\begin{theorem}[$Q_{\mathrm{eff}}$ from S-matrix unitarity]
\label{P17:thm:Qeff_derived}
\begin{equation}
  \boxed{Q_{\mathrm{eff}}
  \;=\; 2\pi\,N_3\,S^{(1)}_{0,0}
  \;=\; \frac{2\pi}{\sqrt{k+2}}\bigg|_{k=3}
  \;=\; \frac{2\pi}{\sqrt{5}},}
  \label{P17:eq:Qeff_theorem}
\end{equation}
where $N_3=\sqrt{2/(k+2)}$ and $S^{(1)}_{0,0}=1/\!\sqrt{2}$ are both
uniquely fixed by unitarity of the respective S-matrices.
\end{theorem}

\begin{proof}
\textbf{Step~1.} Unitarity of $S^{(3)}$ at $j=0$:
$N_3^2\!\sum_{j'=0}^{3/2}\sin^2\!\bigl((2j'{+}1)\pi/5\bigr)=1$.
The identity
$\sum_{j=0}^{k/2}\sin^2\!\bigl((2j{+}1)\pi/(k{+}2)\bigr)=(k{+}2)/2$
gives $N_3^2\cdot\tfrac52=1$, so $N_3=\sqrt{2/5}$.

\textbf{Step~2.} The same identity at $k=1$:
$N_1=\sqrt{2/3}$ and $S^{(1)}_{0,0}=N_1\sin(\pi/3)=1/\!\sqrt{2}$.

\textbf{Step~3.} The effective coupling is the product of the modular
period $2\pi$, the numerator normalisation $N_3$ inherited from $S^{(3)}$,
and the denominator vacuum amplitude $S^{(1)}_{0,0}$ (the coset
quotients out the $\SU(2)_1$ contribution):
\[
  Q_{\mathrm{eff}}
  = 2\pi\cdot N_3\cdot S^{(1)}_{0,0}
  = 2\pi\cdot\sqrt{\tfrac{2}{5}}\cdot\tfrac{1}{\sqrt{2}}
  = \frac{2\pi}{\sqrt{5}}.
\]
No free parameter enters.
\qed
\end{proof}

\begin{remark}[$S$-matrix estimate of $\theta_{12}$]
\label{P17:rem:Smat_phi}
leading-order S-matrix ansatz: using the row of $S^{(3)}$
corresponding to $j=\tfrac{1}{2}$ (the electron-generation
primary), the ratio
$|S_{1/2,1/2}^{(3)}|/|S_{1/2,0}^{(3)}|=\sin(2\pi/5)/\sin(\pi/5)=\vp$
gives a solar angle $\theta_{12}\approx31.7^\circ$, compared to
the experimental $33.41^\circ$ ($1.7^\circ$ off).
The reactor angle $\theta_{13}$ and atmospheric angle $\theta_{23}$
require the full branching-coefficient analysis and are not
predicted at this order.
\end{remark}

\begin{remark}[$m_{\nu_1}$ is strongly suppressed]
\label{P17:rem:m1_suppressed}
The inner gap $\Delta T_{\mathrm{in}}=2/5=4\,\Delta T_{\mathrm{out}}$
gives $m_{\nu_2}/m_{\nu_1}=\exp(2\pi Q_{\mathrm{eff}}\cdot2/5)
=(m_{\nu_3}/m_{\nu_2})^4\approx1167$.  Since
$m_{\nu_2}\approx\sqrt{\Delta m^2_{21}}\approx8.7\,\mathrm{meV}$,
\begin{equation}
  m_{\nu_1} \;\approx\; \frac{8.7\,\mathrm{meV}}{1167}
             \;\approx\; 0.007\,\mathrm{meV},
  \label{P17:eq:m1_pred}
\end{equation}
strongly suppressed relative to $m_{\nu_2}$.  The spectrum is
thus \emph{strictly hierarchical} in the normal ordering.
This is a falsifiable prediction: if absolute neutrino mass
measurements (e.g.\ KATRIN, Project~8, or future experiments)
establish $m_{\nu_1}\gtrsim0.1\,\mathrm{meV}$, the coset
generation assignment in its current form would be disfavoured.
\end{remark}

\subsection{The trefoil $T_{2,3}$ as the geometric source of the coset}
\label{P17:sec:coset_trefoil}

The $\SU(2)_3/\SU(2)_1$ coset does not arise from two independently
chosen WZW models: every parameter of the coset is determined by the
winding numbers $(p,q)=(2,3)$ of the trefoil $T_{2,3}$ itself.

\begin{proposition}[$(p,q)$ determines the coset]
\label{P17:prop:pq_coset}
For the trefoil $T_{p,q}=T_{2,3}$ with $p=2$, $q=3$, the following
identities hold exactly:
\begin{align}
  k_{\mathrm{neutrino}} &= q - p = 1, &
  k_{\mathrm{lepton}} &= q = 3, &
  k_{\mathrm{lepton}}+2 &= p+q = 5, \label{P17:eq:level_from_pq}\\
  Q_1 &= 2(q-p+2) = 6, &
  Q_3 &= 2(p+q) = 10, &
  Q_{\mathrm{coset}} &= \mathrm{lcm}(Q_1,Q_3) = 30,
  \label{P17:eq:Qgroup_from_pq}\\
  Q_{\mathrm{eff}} &= \frac{2\pi}{\sqrt{p+q}} = \frac{2\pi}{\sqrt{5}}, &
  \Delta T_{\mathrm{out}} &= \frac{2q-p}{8(p+q)} = \frac{1}{10}, &
  \frac{\Delta T_{\mathrm{in}}}{\Delta T_{\mathrm{out}}} &= 2q-p = 4.
  \label{P17:eq:gaps_from_pq}
\end{align}
The crossing number of $T_{2,3}$ is $q(p-1)=3$ and equals
$k_{\mathrm{lepton}}$; the writhe is $q=3=k_{\mathrm{lepton}}$.
\end{proposition}

\begin{proof}
Equations~\eqref{P17:eq:level_from_pq} are exact substitutions of
$p=2$, $q=3$. The Q-group formulas in~\eqref{P17:eq:Qgroup_from_pq}
follow from $Q_k=2(k+2)$ (the standard T-matrix closure for $\SU(2)_k$
at $j=\frac{1}{2}$) and $Q_{\mathrm{coset}}=\mathrm{lcm}(Q_1,Q_3)=
\mathrm{lcm}(6,10)=30$ (Theorem~\ref{P17:thm:coset_bosonic}).  The
outer gap formula in~\eqref{P17:eq:gaps_from_pq} follows by direct
computation from Proposition~\ref{P17:prop:coset_T}:
$\Delta T_{\mathrm{out}}=T^{\mathrm{coset}}_{(1/2,1/2)}
-T^{\mathrm{coset}}_{(0,0)}=(-2/15)-(-1/30)=-(3/30)=1/10$
(in absolute value),
and $(2q-p)/(8(p+q))=4/40=1/10$.  The gap ratio $4=k_1+k_3$
follows from Proposition~\ref{P17:prop:coset_gaps}.  The crossing
number $q(p-1)=3$ is the standard formula for $T_{p,q}$ with $p<q$.
\qed
\end{proof}

\begin{remark}[Coset framing independence via the writhe]
\label{P17:rem:framing}
The trefoil $T_{2,3}$ has writhe $w=q=3$ (the $\QH=3$ baryon sector's
own topological charge).  The neutrino sector ($\QH=1$, proved unknot
with writhe $0$ in Proposition~\ref{P17:prop:qh1_unknot}) has $w=0$.  The
coset T-matrix $T^{\mathrm{coset}}=T^{(3)}-T^{(1)}$ subtracts the
two sectors' T-matrices, making the coset automatically \emph{framing-independent}: the trefoil's writhe cancels between the
two sector contributions.  This is a structural guarantee, not a
tuning: the coset is the correct mathematical object for the
inter-sector (neutrino--charged-lepton) relationship precisely
because the difference of T-matrices removes the Q$_H$=3 sector's
self-framing.
\end{remark}

\begin{remark}[Kauffman bracket interpretation of the four coset fields]
\label{P17:rem:kauffman}
The trefoil $T_{2,3}$ has three crossings, giving $2^3=8$ Kauffman
smoothing states (A- or B-smoothing at each crossing).  Grouping by
the number $n$ of resolved crossings gives four sectors with
multiplicities $\binom{3}{n}$:
\begin{center}
\renewcommand{\arraystretch}{1.3}
\begin{tabular}{ccccc}
\toprule
resolved $n$ & $2\Lambda=n$ & $\binom{3}{n}$ & $T^{\mathrm{coset}}$
  & mass order \\
\midrule
$0$ & $0$  & $1$ & $-1/30$  & vacuum (ground state) \\
$1$ & $1$  & $3$ & $-2/15$  & \textbf{heaviest} ($\nu_3$) \\
$2$ & $2$  & $3$ & $11/30$  & second lightest ($\nu_2$) \\
$3$ & $3$  & $1$ & $7/15$   & lightest ($\nu_1$) \\
\bottomrule
\end{tabular}
\end{center}
The four coset primaries~(Proposition~\ref{P17:prop:coset_T}) are
identified with these four sectors via $2\Lambda=$ number of
resolved crossings, where $\Lambda\in\{0,\tfrac12,1,\tfrac32\}$
is the $\SU(2)_3$ spin label.  The mass ordering is \emph{inverted}
relative to the naive expectation: the state with one crossing
resolved is heavier than the fully-bound state ($n=0$).  This
inversion is the analogue of string tension: partially liberating
one quark from the baryon (resolving one crossing) stretches the
remaining two-string configuration to maximum tension, storing
more energy than the fully-confined ground state.  The state $n=1$
(one crossing resolved, maximum tension) is the heaviest; $n=3$
(all crossings resolved, no tension) is the lightest.

The $\mathbb{Z}_3$ symmetry of the trefoil's three crossings is the
same $\mathbb{Z}_3$ as the 3-state Potts model symmetry of the coset
$M(5,6)$: the geometric symmetry of the quark positions \emph{is}
the algebraic symmetry of the inter-sector mixing theory.
\end{remark}

\begin{theorem}[$Q_{\mathrm{coset}}=h(E_8)$: the Brieskorn sphere as common source]
\label{P17:thm:brieskorn}
$Q_{\mathrm{coset}}=h(E_8)=30$.  Both values are determined by the
Brieskorn sphere $\Sigma(2,3,5)=\{z^2+w^3+v^5=0\}\cap S^5\subset\mathbb{C}^3$,
which is the single geometric object encoding both $T_{2,3}$ and $2I$
simultaneously.  More precisely:
\begin{enumerate}[label=\emph{(\roman*)}]
\item $Q_{\mathrm{coset}}=\bigl(-1+\tfrac{1}{p}+\tfrac{1}{q}
  +\tfrac{1}{p+q}\bigr)^{-1}=\bigl(-1+\tfrac{1}{2}+\tfrac{1}{3}+\tfrac{1}{5}\bigr)^{-1}=30$,
  where $(p,q)=(2,3)$ are the winding numbers of the trefoil
  (Proposition~\ref{P17:prop:pq_coset}).
\item $h(E_8)=|\twoI|/4=|\pi_1(\Sigma(2,3,5))|/4=120/4=30$, where
  $\pi_1(\Sigma(2,3,5))=\twoI$ is the binary icosahedral group and
  $h(E_8)$ is its Coxeter number.
\item These are equal because the arithmetic identity $8q=(p+q)^2-1$
  holds at $(p,q)=(2,3)$~($8\times3=24=(5)^2-1$~$\checkmark$), which
  gives $(p+q)\bigl((p+q)^2-1\bigr)/4=2q(p+q)$.  Among all torus
  knots $T_{p,q}$ with $\gcd(p,q)=1$ and $p\geq2$, $T_{2,3}$ is
  the unique knot satisfying this identity.
\end{enumerate}
\end{theorem}

\begin{proof}
Part~(i) is the orbifold Euler characteristic formula for the base
orbifold $S^2(p,q,p+q)$ of the Seifert fibration of $\Sigma(p,q,p+q)$:
$\chi=2-3+\tfrac{1}{p}+\tfrac{1}{q}+\tfrac{1}{p+q}=-1+\tfrac{1}{2}+\tfrac{1}{3}+\tfrac{1}{5}=\tfrac{1}{30}$.
Then $Q_{\mathrm{coset}}=1/\chi=30$, consistent with Proposition~\ref{P17:prop:pq_coset}
($Q_{\mathrm{coset}}=2q(p+q)=30$).

Part~(ii): The classical result $\pi_1(\Sigma(2,3,5))=\twoI$ is due to
Brieskorn~\cite{Brieskorn1966}.  The identity $|\twoI|=|\mathrm{SL}(2,5)|
=5(5^2-1)=120$ gives $h(E_8)=|\twoI|/4=30$.

Part~(iii): The two expressions obtained independently in~(i) and~(ii)
are equal at $(p,q)=(2,3)$ because $2q(p+q)=(p+q)\bigl((p+q)^2-1\bigr)/4$
is equivalent (on dividing by $p+q>0$) to $8q=(p+q)^2-1$, which
evaluates to $24=24$ at $(p,q)=(2,3)$.  Uniqueness: a search over all
torus knots $T_{p,q}$ with $p\geq2$, $q>p$, $\gcd(p,q)=1$ finds no
other solution; $T_{2,3}$ is the only knot in this class for which
the orbifold formula and $|\twoI|/4$ agree.

Combining: $Q_{\mathrm{coset}}=1/\chi(\Sigma(2,3,5))
=|\pi_1(\Sigma(2,3,5))|/4=h(E_8)$.
\qed
\end{proof}

\begin{remark}[The three sectors inside $\Sigma(2,3,5)$]
\label{P17:rem:brieskorn_sectors}
The Brieskorn sphere $\Sigma(p,q,p+q)=\Sigma(2,3,5)$ encodes all
three physical sectors simultaneously:
\begin{itemize}[nosep]
\item The variety $\{z^2+w^3=0\}$ links to the trefoil $T_{2,3}$:
  the $\QH=3$ \emph{baryon sector}.
\item $\pi_1(\Sigma(2,3,5))=\twoI$: the McKay group of the $\QH=2$
  \emph{charged-lepton sector} (Paper~XVI).
\item The three Seifert invariants $(p,q,p+q)=(2,3,5)$ are
  $(k_\nu+1,\,k_\ell,\,k_\ell+2)$, encoding the WZW levels of the
  $\QH=1$ \emph{neutrino sector} ($k_\nu=1$) and $\QH=2$ sector
  ($k_\ell=3$, denominator $k_\ell+2=5$).
\end{itemize}
The orbifold Euler characteristic $\chi=-1+\tfrac12+\tfrac13+\tfrac15=\tfrac{1}{30}$
then simultaneously gives $Q_{\mathrm{coset}}=30=h(E_8)=|\twoI|/4$.
\end{remark}

\begin{remark}[$Q_{\mathrm{coset}}=h(E_8)$ and the strange quark]
\label{P17:rem:strange_quark}
Theorem~\ref{P17:thm:brieskorn} resolves the open question of whether the
identity $Q_{\mathrm{coset}}=h(E_8)=30$ is structural or coincidental:
it is structural, both values being $|\pi_1(\Sigma(2,3,5))|/4$.

The physical consequence concerns the strange quark.
The E$_8$ Coxeter spectrum gives the strange quark the independently-derived
T-matrix coincidence $T^{(m=13)}=1/8=T_\mu$ of Paper~XVI
(Remark~\ref{P17:rem:strange_quark}), confirmed to $0.09\%$ after one-loop
RG running), where $m=13$ is the 4th exponent of $E_8$. Now 
$h(E_8)=|\twoI|/4$ is the same integer as $Q_{\mathrm{coset}}=
|\pi_1(\Sigma(2,3,5))|/4$, and the crossing-resolution picture
(Remark~\ref{P17:rem:kauffman}) provides the physical link:

In the baryon sector, the three trefoil crossings are the three quark
positions.  One crossing resolving corresponds topologically to a transition
$T_{2,3}\to$ (Hopf link), i.e.\ $\QH=3\to\QH=2$.  The strange quark,
being the lightest quark carrying a quantum number (strangeness) absent
from the nucleon ground state, is the natural candidate for the quark at
the $\QH=3\to\QH=2$ threshold crossing.  The Brieskorn sphere
$\Sigma(2,3,5)$ is the geometric object that bridges the trefoil (source
of the crossing) and $\twoI$ (source of the E$_8$ exponent $m=13$): both
arise as different projections of the same $\Sigma(2,3,5)$, so the
$Q_{\mathrm{coset}}=h(E_8)$ identity is the algebraic shadow of this
common geometric origin. Liberating the strange quark is precisely the transition from
the SU(2)-isospin symmetric world ($u,d$ quarks) to the SU(3)-flavour-broken
world ($u,d,s$), which is topologically the transition from one crossing
resolved to two.

\end{remark}

%══════════════════════════════════════════════════════════════════════════════
\section{Open Problems}
\label{P17:sec:open}
%══════════════════════════════════════════════════════════════════════════════

Section~\ref{P17:sec:coset} resolves the question of which inter-sector
structure encodes neutrino generation labels (the coset
$\SU(2)_3/\SU(2)_1$), provides a $0.9\%$-accurate prediction for
$m_{\nu_3}/m_{\nu_2}$, and derives $Q_{\mathrm{eff}}=2\pi/\sqrt{k+2}$
from S-matrix unitarity
(Proposition~\ref{P17:prop:Qeff_unique}, Theorem~\ref{P17:thm:Qeff_derived}).
Section~\ref{P17:sec:nu_cmb_matching} completes the $c=1$ thermodynamic
matching, deriving $\mathcal{C}^{(\nu)}=0$ exactly and
predicting $m_{\nu_1}\approx42\,\mathrm{meV}$ with $n_\nu=42.39$.
Section~\ref{P17:sec:k1_dynamics} establishes $k=1$ via the Wess--Zumino
descent and proves Majorana statistics (Corollary~\ref{P17:cor:majorana}).
Section~\ref{P17:sec:pmns_branching} determines the primary-level PMNS
branching structure exactly, deriving the leading-order permutation
matrix (Theorem~\ref{P17:thm:pmns_permutation}), computing the first
subleading branching coefficients, and proving the electron's exact
decoupling from all neutrino mass eigenstates at every order.
The remaining open problem is the full $q$-series resummation that
converts the branching functions into physical mixing angles.

\medskip\noindent
There are no enumerated open problems remaining in this paper beyond
the PMNS resummation, which is developed to partial results in
Section~\ref{P17:sec:pmns_branching} and will be completed in Papers~XIX--XX.

%══════════════════════════════════════════════════════════════════════════════
\section{Summary}
\label{P17:sec:summary}
%══════════════════════════════════════════════════════════════════════════════

This paper characterises the $\QH=1$ sector of the density-feedback
Faddeev--Niemi Hopfion and identifies it as the neutrino sector. The
structural results are:

\begin{enumerate}[noitemsep]
\item The $\QH=1$ Hopfion's preimage of any regular value is a single
  unknotted circle linking any reference fibre with linking number~$1$
  (Proposition~\ref{P17:prop:qh1_unknot}). Writhe is zero; no self-linking
  or crossing structure is present.
\item The McKay-correspondence group for $\QH=1$ is
  $\Ztwo=\{\pm1\}$, with McKay chain
  $\Ztwo\leftrightarrow\widehat{A}_1\leftrightarrow\SU(2)_1$, placing
  the sector in the A-type (cyclic, abelian) series of the ADE
  classification and distinguishing it from the E-type exceptional
  groups that organise $\QH=2$ and $\QH=3$
  (Proposition~\ref{P17:prop:qh1_Z2}).
\item Colour exclusion for $\QH=1$ holds by a simple coprimality
  argument: $\gcd(|\Ztwo|,|\Zthree|)=\gcd(2,3)=1$, so the colour
  $\Zthree$ has no non-trivial action on $\Ztwo$-representations
  (Corollary~\ref{P17:cor:qh1_colour_free}). This strengthens
  Paper~XVI's Theorem~8.1%thm:colour_exclusion
  , which required the detailed structure of $\mathrm{Out}(\twoI)$ and
  $N_{\twoI}(\twoT)$.
\item The $\QH=1$ sector has no crossing regions, no Y-junction, and
  zero string tension $\sigma=0$ by direct topology
  (Proposition~\ref{P17:prop:qh1_noconfinement}), making it the
  maximally free sector in the confinement hierarchy
  $\QH=1\subsetneq\QH=2\subsetneq\QH=3$.
\item The T-matrix closure condition for $\SU(2)_1$ closes at
  $Q_{\mathrm{group}}=6$ (Proposition~\ref{P17:prop:qh1_Qgroup}), the
  smallest Q-group in the hierarchy $\{3,6,10\}$ for
  $\{\QH=3,\QH=1,\QH=2\}$. The three neutrino-generation T-matrix
  phases are $T_g^{(\nu)}=5/24,\,1/8,\,3/40$ for $g=1,2,3$, following
  the same level-ladder formula as the charged-lepton sector.
\item The quantum dimension $\dimq{1/2}=\vp$ at level $k=3$ is exact
  and equal to the $\QH=2$ sector's $\dimq{1}=\vp$
  (Proposition~\ref{P17:prop:qh1_dimq}). Both unconfined lepton sectors
  share quantum dimension $\vp$; the confined baryon sector has
  $\dimq{3/2}=1$ (the simple-current condition).
\item The Wess--Zumino descent (Theorem~\ref{P17:thm:kone_forced}) forces
  the $\QH=1$ condensate's own WZW level to $k=1$, not the ambient
  $k=3$: the McKay relation $Q_{\mathrm{group}}^{(\nu)}=2(k+2)=6$ gives
  $k=1$ uniquely.  The sector-specific virial is $V^*=\dim_q(\tfrac12;k{=}1)=1$
  (not $\vp$), and the profile ratio changes to
  $\Jfour^{(\nu)}/\Ja^{(\nu)}=(1+\sqrt{3}/\vp)\cdot2^{1/3}/\vp^6\approx0.145$
  (Proposition~\ref{P17:prop:k1_virial}).  The thin-torus normalization identity at $k=1$
  gives $5\vp\sqrt{2/(\vp+\sqrt{3})}\approx6.25=6\cdot[1+O(R_0^{-2})]$
  (Proposition~\ref{P17:prop:norm_id_k1}), consistent with two-route
  equivalence holding at $k=1$ and the mass prediction
  $m_{\nu_1}\approx42\,\mathrm{meV}$ remaining unchanged
  (Corollary~\ref{P17:cor:two_route_k1}).
\item The $j=\tfrac12$ primary of $\mathrm{SU}(2)_1$ is self-conjugate
  under charge conjugation ($j^*=k-j=\tfrac12=j$) and generates a
  $\Ztwo$ fusion group ($\tfrac12\times\tfrac12=0$): it is a Majorana
  simple current.  The $\QH=1$ condensate therefore predicts Majorana
  neutrinos from the WZW structure alone
  (Corollary~\ref{P17:cor:majorana}).
\item Under Conjecture~\ref{P17:conj:nu_identification} (neutrino
  identification), the exact experimental bound is $n_\nu\geq42.02$
  (Proposition~\ref{P17:prop:nu_tower_bound}), with no integer rounding required
  (Remark~\ref{P17:rem:nconv}).  The $c=1$ thermodynamic
  matching of Section~\ref{P17:sec:nu_cmb_matching} gives the exact
  irrational prediction $n_\nu=42.39$ and $m_{\nu_1}\approx42\,\mathrm{meV}$
  (Theorem~\ref{P17:thm:qh1_tworoute}, Proposition~\ref{P17:prop:CS_spoke},
  Corollary~\ref{P17:cor:qh1_mass}), using $\mathcal{C}^{(\nu)}=0$ proved
  from Paper~IV~\cite{Paper4}.  The prediction satisfies the bound
  with margin~$0.37$; the former ``one-unit gap'' is resolved by
  Paper~VI~\cite{Paper6}'s irrationality theorem for tower levels
  (tower levels are irrational for all lepton generations, so integer
  rounding is not applicable).  The master mass formula
  (Remark~\ref{P17:rem:master_formula}) unifies this result with the
  electron mass under a single expression parametrised by $\TCMB$ alone.
\item The inter-sector coset $\SU(2)_3/\SU(2)_1$ carries four primary
  fields with T-matrix phases
  $\{-\frac{2}{15},-\frac{1}{30},\frac{11}{30},\frac{7}{15}\}$
  (Proposition~\ref{P17:prop:coset_T}). The coset is \emph{bosonic}
  ($Q_{\mathrm{coset}}=30$, $4Q\cdot T^{\mathrm{coset}}\in2\mathbb{Z}$
  for all fields; Theorem~\ref{P17:thm:coset_bosonic}), in contrast to the
  fermionic matter sectors. The four T-matrix values form the pattern
  $\{-4,-1,11,14\}/30$ with gap structure $(1/10,\,2/5,\,1/10)$
  (Proposition~\ref{P17:prop:coset_gaps}).
\item The atmospheric pair $\{(1,0),(\frac{3}{2},\frac{1}{2})\}$,
  separated by the outer gap $\Delta T=1/10$, gives
  $m_{\nu_3}/m_{\nu_2}=\exp(4\pi^2/10\sqrt{5})\approx5.845$ at
  $Q_{\mathrm{eff}}=2\pi/\sqrt{5}$, matching the experimental
  value $5.795$ to $0.9\%$ (Proposition~\ref{P17:prop:nu_mass_ratio}).
  The value $Q_{\mathrm{eff}}=2\pi/\sqrt{k+2}$ is uniquely forced
  by bosonic monodromy (irrationality) and S-matrix unitarity
  (Proposition~\ref{P17:prop:Qeff_unique}, Theorem~\ref{P17:thm:Qeff_derived}).
  The inner gap $\Delta T_{\mathrm{in}}=2/5=4\,\Delta T_{\mathrm{out}}$
  predicts $m_{\nu_1}\approx0.007\,\mathrm{meV}$
  (Remark~\ref{P17:rem:m1_suppressed}).
\item Every parameter of the coset --- level assignments,
  Q-groups, gap structure, effective Q, and framing independence ---
  is determined exactly by the trefoil's winding numbers $(p,q)=(2,3)$:
  $k_{\mathrm{neutrino}}=q-p=1$, $k_{\mathrm{lepton}}=q=3$,
  $Q_1=6$, $Q_3=10$, $Q_{\mathrm{coset}}=\mathrm{lcm}(Q_1,Q_3)=30$,
  $\Delta T_{\mathrm{out}}=(2q-p)/(8(p+q))=1/10$
  (Proposition~\ref{P17:prop:pq_coset}). The trefoil does not merely
  realise the baryon sector; its $(p,q)$ parameters encode the
  inter-sector coset that governs neutrino mass generation.
\item The primary-level branching of $\SU(2)_3\to\SU(2)_1\otimes M(5,6)$
  gives the PMNS matrix \emph{exactly} as the permutation
  (Theorem~\ref{P17:thm:pmns_permutation}, Corollary~\ref{P17:cor:pmns_exact},
  equation~\eqref{P17:eq:PMNS_LO}): $\mu\to\nu_3$, $\tau\to\nu_2$,
  unitarity forces $e\to\nu_1$, with
  $\theta_{23}=45^\circ$, $\theta_{13}=0^\circ$, $\theta_{12}=0^\circ$.
  The permutation is the complete coset prediction, not a leading-order
  approximation: each generation branches into exactly one sector at
  every descendant level, leaving no mechanism for off-diagonal
  corrections within the coset branching.
  The electron decouples from all neutrino mass eigenstates exactly at
  every level (Corollary~\ref{P17:cor:electron_decouple}).
  The experimental values $\theta_{23}\approx48^\circ$,
  $\theta_{13}\approx8.5^\circ$, $\theta_{12}\approx33^\circ$
  (deviations $1/(k+2)$ to $1/(k+2)^2$) arise from RG running
  or condensate-background mixing, not from the coset branching itself
  (Section~\ref{P17:sec:pmns_corrections}).
\item The four coset fields correspond to the four Kauffman
  crossing-resolution sectors of $T_{2,3}$'s three crossings,
  with $2\Lambda=n_{\mathrm{resolved}}$ (Remark~\ref{P17:rem:kauffman}).
  The $n=1$ state is the heaviest ($\nu_3$): partially liberating
  one quark maximises string tension, as in QCD.  The identity
  $Q_{\mathrm{coset}}=h(E_8)=30$ is structural: both equal
  $|\pi_1(\Sigma(2,3,5))|/4$, where the Brieskorn sphere
  $\Sigma(2,3,5)$ encodes $T_{2,3}$ and $\twoI$ as projections of
  the same object (Theorem~\ref{P17:thm:brieskorn}), connecting the
  strange quark's T-phase coincidence $T_s=T_\mu$ of Paper~XVI
  to the coset geometry directly.
\end{enumerate}

The sector hierarchy is now complete through $\QH=3$:
\begin{center}
\renewcommand{\arraystretch}{1.35}
\begin{tabular}{lllllll}
\toprule
$\QH$ & Particle & Group & McKay & WZW & $Q_{\mathrm{gr}}$ & $\dimq{}$ \\
\midrule
$0$ & Vacuum & $\{1\}$ & --- & --- & --- & $1$ \\
$1$ & Neutrino & $\Ztwo$ & $\widehat{A}_1$ & $\SU(2)_1$ & $6$ & $\vp$ \\
$2$ & Charged lepton & $\twoI$ & $\widehat{E}_8$ & $\SU(2)_3$ & $10$ & $\vp$ \\
$3$ & Baryon & $\twoT$ & $\widehat{E}_6$ & $(E_6)_1$ & $3$ & $1$ \\
\bottomrule
\end{tabular}
\end{center}
The jump from A-type ($\QH=1$) to E-type ($\QH=2,3$) at $\QH=2$
corresponds precisely to the acquisition of electroweak charge, while
the E-type to E-type transition $\twoI\to\twoT$ at $\QH=3$ corresponds
to the emergence of colour charge (Paper~XVI,
Theorem~8.1%thm:colour_exclusion
).

%══════════════════════════════════════════════════════════════════════════════
\begin{thebibliography}{99}

\bibitem{Paper1}
F.~Manfredi,
\textit{The Density-Feedback Faddeev--Niemi Hopfion:
  Exact Algebraic Predictions for Standard-Model Parameters},
Zenodo (2026).
\href{https://doi.org/10.5281/zenodo.19342027}{doi:10.5281/zenodo.19342027}

\bibitem{Paper3}
F.~Manfredi,
\textit{Two Constants from One Knot: The Fine Structure Constant and
  the Linking Scale as Exact Predictions of the Icosahedral WZW Condensate},
Zenodo (2026).
\href{https://doi.org/10.5281/zenodo.19478629}{doi:10.5281/zenodo.19478629}

\bibitem{Paper4}
F.~Manfredi,
\textit{The Hopf Spoke, WZW Fermion Mass Renormalization,
  and Three- and Four-Term Formulas for the Fine Structure Constant},
Zenodo (2026).
\href{https://doi.org/10.5281/zenodo.19504729}{doi:10.5281/zenodo.19504729}

\bibitem{Paper5}
F.~Manfredi,
\textit{Colour Confinement, the QCD Scale, and Hadronic Structure
  from the Density-Feedback Hopfion},
Zenodo (2026).
\href{https://doi.org/10.5281/zenodo.19638857}{doi:10.5281/zenodo.19638857}

\bibitem{Paper6}
F.~Manfredi,
\textit{The Golden-Ratio Tower: Fermion Scale Hierarchy from the Hopfion RG Fixed Point},
Zenodo (2026).
\href{https://doi.org/10.5281/zenodo.19646945}{doi:10.5281/zenodo.19646945}

\bibitem{Paper12}
F.~Manfredi,
\textit{The Profile Normalisation Conjecture ---
  Proof via CMB Energy Identity and Two-Route Equivalence},
Zenodo (2026).
\href{https://doi.org/10.5281/zenodo.20210460}{doi:10.5281/zenodo.20210460}

\bibitem{Paper14}
F.~Manfredi,
\textit{The Thick-Torus Profile Correction to the Density-Feedback Hopfion},
Zenodo (2026).
\href{https://doi.org/10.5281/zenodo.20479790}{doi:10.5281/zenodo.20479790}

\bibitem{Paper15}
F.~Manfredi,
\textit{The $Q_H=3$ Sector of the Density-Feedback Hopfion},
Zenodo (2026).
\href{https://doi.org/10.5281/zenodo.20691001}{doi:10.5281/zenodo.20691001}

\bibitem{Paper16}
F.~Manfredi,
\textit{Quark Generation Masses in the Density-Feedback Hopfion:
A Survey of Ruled-Out Mechanisms and the $Q_H=3$ Gradient-Flow Landscape},
Preprint (2026).
\href{https://doi.org/10.5281/zenodo.21012650}{doi:10.5281/zenodo.21012650}

\bibitem{Paper18}
F.~Manfredi,
\textit{Preimage Topology and Confinement: The $Q_H=3$ Sector of the Density-Feedback Hopfion},
Zenodo (2026).
\href{https://doi.org/10.5281/zenodo.21047750}{doi:10.5281/zenodo.21047750}

\bibitem{Paper19}
F.~Manfredi,
\textit{Dynamical Origin of Quark Masses in the Density-Feedback Hopfion},
Zenodo (2026).
\href{https://doi.org/10.5281/zenodo.21225452}{doi:10.5281/zenodo.21225452}

\bibitem{McKay1980}
J.~McKay,
\textit{Graphs, singularities, and finite groups},
Proc.\ Symp.\ Pure Math.\ \textbf{37}, 183--186 (1980).
\href{https://doi.org/10.1090/pspum/037}{doi:10.1090/pspum/037}

\bibitem{Faddeev1997}
L.~D. Faddeev and A.~J. Niemi,
\textit{Stable knot-like structures in classical field theory},
Nature \textbf{387}, 58 (1997).
\href{https://doi.org/10.1038/387058a0}{doi:10.1038/387058a0}

\bibitem{PDG2024}
S.~Navas et al.\ (Particle Data Group),
\textit{Review of Particle Physics},
Phys.\ Rev.\ D \textbf{110}, 030001 (2024).
\href{https://doi.org/10.1103/PhysRevD.110.030001}{doi:10.1103/PhysRevD.110.030001}

\bibitem{Kac1990}
V.~G.~Kac,
\textit{Infinite Dimensional Lie Algebras}, 3rd ed.,
Cambridge University Press (1990).
ISBN 978-0-521-46693-6

\bibitem{Brieskorn1966}
E.~Brieskorn,
\textit{Beispiele zur Differentialtopologie von Singularit\"aten},
Inventiones Mathematicae \textbf{2}, 1--14 (1966).
\href{https://doi.org/10.1007/BF01403388}{doi:10.1007/BF01403388}

\bibitem{Asner2015}
Project~8 Collaboration (D.~M.~Asner et al.),
\textit{Single-electron detection and spectroscopy via relativistic
  cyclotron radiation},
Phys.\ Rev.\ Lett.\ \textbf{114}, 162501 (2015);
see also Project~8 Collaboration (A.~Ashtari~Esfahani et al.),
\textit{The Project~8 neutrino mass experiment},
arXiv:2203.07349 (2022), target sensitivity $m_\beta\geq 40$~meV/$c^2$
at $90\%$~CL.
\href{https://doi.org/10.1103/PhysRevLett.114.162501}{doi:10.1103/PhysRevLett.114.162501}

\bibitem{Planck2018}
Planck Collaboration,
\textit{Planck 2018 results. VI. Cosmological parameters},
A\&A \textbf{641}, A6 (2020).
\href{https://doi.org/10.1051/0004-6361/201833910}{doi:10.1051/0004-6361/201833910}

\end{thebibliography}
\end{document}